% !TeX program = xelatex
\documentclass[journal,onecolumn]{IEEEtran}
\usepackage{iftex}
\ifPDFTeX
  \errmessage{This project requires XeLaTeX. In Overleaf: Menu -> Compiler -> XeLaTeX}
\fi

\usepackage{geometry}
\geometry{margin=2.3cm}
\usepackage{amsmath,amssymb,amsthm,mathtools}
\usepackage[strings]{underscore}
\usepackage{booktabs}
\usepackage{longtable}
\usepackage{graphicx}
\graphicspath{{figures/}}
\usepackage{titlesec}
\usepackage{hyperref}
\hypersetup{hidelinks}
\usepackage{placeins}
\usepackage{float} % provides [H] float placement

\usepackage{fontspec}
\usepackage{xeCJK}
\setmainfont{Times New Roman}
\setCJKmainfont{FandolSong-Regular}
\linespread{1.2}
\renewcommand{\thesection}{\arabic{section}}
\renewcommand{\thesubsection}{\thesection.\arabic{subsection}}
\renewcommand{\thesubsubsection}{\thesubsection.\arabic{subsubsection}}
\titleformat{\section}{\normalfont\Large\bfseries}{\thesection}{0.6em}{}
\titleformat{\subsection}{\normalfont\large\bfseries}{\thesubsection}{0.6em}{}
\titleformat{\subsubsection}{\normalfont\normalsize\bfseries}{\thesubsubsection}{0.6em}{}

\title{论文提纲v2}
\author{匿名作者}
\date{}

\begin{document}
\maketitle

\begin{abstract}
本文为论文提纲，面向可审计半自动全景布局标注流程，覆盖研究问题、方法、实验设置、报告口径与预期贡献。
\end{abstract}

\begin{IEEEkeywords}
%Semi-automatic annotation, Reliability, Routing, OOD, Auditable protocol
\end{IEEEkeywords}

\tableofcontents

\section{研究问题}

本研究关注全景室内场景中半自动角点标注的质量保障问题。在无法对所有样本进行全量精标的现实约束下，如何系统性地评估半自动标注的效率、质量与可靠性，并利用场景特异可靠度支持任务路由。

技术主线：采用可审计三段式规程 Pilot→PreScreen→Main；主实验前锁定阈值与规则，包括 MDE、$w_u$ 映射、关键超参数与扰动算子库；在 non-IID/OOD 条件下以标注前可得的代理信号（如 $d_t$ 与 $g_t$）驱动首轮风险分层与应激触发，并以标注后 meta-label 共识作诊断与审计；输出画像、矩阵与反例库等过程证据。
\subsection{研究问题 RQ}

本文围绕三个核心研究问题展开，分别对应效率、质量与路由三个维度。

\textbf{RQ1 效率评估}：半自动初始化是否降低真实活跃标注时间？下降比例是多少？
\\ 主终点：Manual vs Semi-Auto 的 \texttt{active\_time} 中位数差异（秒）及 95\% CI。 \texttt{active\_time} 的操作定义、采集实现与聚合口径见第~\ref{sec:method-active-time}节。
\\ 风险控制：分析 \texttt{active\_time} 与 $\mathrm{IoU}_{\mathrm{edit}}$ 的关系，排除时间缩短仅源于随意修改。

\textbf{RQ2 质量与可靠性}：在部分复核约束下，半自动流程是否改变标注一致性与可靠性分布，能否系统性暴露反例类型？
\\ 主终点：在预注册的同图配对子集（$N_{img}=25$）上比较 $\mathrm{IAA}_t$ 的变化，效应量取 $\mathrm{median}(\Delta_t)$，并使用配对置换检验与配对 bootstrap 给出 95\% CI；K-S 检验仅作为附录敏感性分析，不作为主终点。
\\ 反例类型：Type 1 低一致性；Type 2 异常编辑；Type 3 几何合法性失败（定义见第~\ref{sec:method-counterexample}节）。
\\ 可复现筛选规则：按第~\ref{sec:method-counterexample}节规则自动筛选候选，最终反例需专家复核与记录原因。

\textbf{RQ2b 加权共识}：加权共识作为干预手段，是否能在不引入循环论证的前提下提升共识稳定性？
\\ 评价口径：主评价使用独立参考子集或未加权一致性指标；加权共识的适用前提、锁定流程与报告方式见第~\ref{sec:method-weighted-consensus}节。

\textbf{RQ3 路由与优化}：在固定标注预算下，能否通过预筛选、工人画像与场景/分布感知路由，在 non-IID 与 OOD 场景中维持一致性并改善预算效率？同时，对门控失败与字段缺失等安全性事件保持可追溯披露，避免以静默过滤换取表面指标。
\\ 场景定义：以 \texttt{difficulty} 与 \texttt{model\_issue} 的共识标签作为场景代理。全部标签进入审计集合并在报告中给出频率、共现与反例分布。为避免小样本下过度细分导致数据稀疏，路由仅在不超过 $S_{\mathrm{core,max}}=4$ 个核心场景上启用，并必须披露启用率与退化率。
\\ 场景特异可靠度：在上述有限场景集合上估计标注者的场景内一致性中位数，并通过启用门槛、启用率/退化率披露与保守 LCB 排序保证稳健性；固定系数收缩仅作为附录敏感性分析报告（不作为主实现）。定义见第~\ref{sec:method-scene-reliability}节。
路由策略：Random；Global（仅用全局 $r_u$）；Full（PreScreen + 场景分层 + OOD 触发），定义与应激情景见第~\ref{sec:exp-iid-stress}节。
\\ 评估协议：离线模拟，在每任务已有的 k 份真实标注中，模拟策略分配任务并取实际提交标注作为输出，以避免因使用真实标签导致的信息泄露与乐观估计。
\\ 指标（主）：一致性代理提升、预算效率（序贯停止的 $\bar{k}_{\mathrm{used}}$ 或等价预算-覆盖关系）。
\\ 安全性/数据可用性审计（非主终点）：门控失败率与失败原因分布、关键字段缺失率；用于证明不静默过滤与口径 T/I/M 的可追溯性，不作为“路由收益”的主证据。



\section{方法}

本章建立改动幅度、质量、可靠性三类指标体系，并明确各指标的适用范围与边界。核心原则如下：
\begin{enumerate}
    \item 改动幅度反映工作量而非正确性；
    \item 质量与可靠性分开统计；
    \item 所有剔除与失败均记录原因以保证可追溯性。
\end{enumerate}

\subsection{总体框架：预筛选、画像、审计、路由}

本研究采用可复现的预筛选与工人画像框架。考虑到小样本条件下边做边收敛的强化学习闭环缺乏可审计性，本文转而强调过程性证据与规程透明性。具体采用 Pilot→PreScreen→Main Experiment 的三段式流程，但锁定并非一次性完成，而是按证据强度分阶段进行：Pilot 阶段锁定协议骨架与不可变边界，PreScreen 阶段锁定 admission 与加权相关阈值，Calibration 第一轮后再锁定依赖场景覆盖的 scene-routing 参数，从而避免在尚无 Calibration 证据时提前决定场景阈值。

整体流程分为四个阶段：
\begin{enumerate}
    \item 预筛选（PreScreen）：用小规模盲测任务集建立工人基线，剔除明显不合格或不稳定的工人，形成可用工人池。
    \item 校准（Calibration）：在 \texttt{Calibration\_manual} 上估计全局可靠度 $r_u$ 与场景特异可靠度 $r_u^{\left(s\right)}$，并输出置信区间。
    \item 在线路由（Routing）：结合标注前可得代理信号（如 $d_t$ 与 $g_t$）与标注后场景信号（meta-label 共识），在预算约束下将高风险任务优先分配给更稳定的工人，并采用序贯冗余与停止准则。
    \item 审计输出（Audit）：输出工人画像图（可靠度主轴 vs 离散风险层级）与类型分布、Worker×Scene 矩阵、门控失败原因分布、反例案例库，形成可追溯证据链。
\end{enumerate}




\subsection{标注对象与元信息}

标注以角点为主，不强制要求画墙面。每次提交时需同步记录三类信息：Scope、Difficulty 与 Model Issue，这三类元标签作为标注决策的副产品，同时服务于后续的场景分层、工人画像与路由策略。
\subsubsection{Scope 单选，必填，决定是否进入主指标}
\begin{itemize}
    \item In-scope：仅标注相机所在房间
    \item 规则：门框与墙垛清晰时，边界止于门框与墙垛处，不跨门洞
    \item 包含：天花内凹与格栅不影响外边界
\end{itemize}
Out-of-scope（共 4 类）：
\begin{itemize}
    \item 几何假设不成立：非 Manhattan 或弧形墙
    \item 边界不可判定：开放式或无可靠停止点
    \item 错层或多平面：楼梯井或分层天花
    \item 证据不足：遮挡或画质差导致无法判断
\end{itemize}
    
统计处理：标记为 In-scope 的样本进入主指标口径 I；标记为 OOS 的样本单独报告其比例、子类分布与判定一致性；当同一图像存在标注者分别选择 In-scope 与 OOS 的混合判定时，进行人工复核。
混合判定的操作流程与裁决规则见第~\ref{sec:method-oos}节。
\subsubsection{Difficulty 多选}
用户界面要求：Difficulty 为多选列表。为避免“空值既可能表示无困难、也可能表示漏填”的不可区分性，界面要求每次提交\textbf{至少选择 1 项}，并施加互斥约束：
\begin{itemize}
    \item 若选择 \texttt{trivial}（表示无明显困难），则\textbf{不得}同时选择任何其他困难标签；
    \item 若存在任意具体困难因素（如 occlusion/seam 等），则\textbf{不得}选择 \texttt{trivial}。
\end{itemize}
若出现空选、或出现 \texttt{trivial} 与其他困难标签共存的冲突选择，均视为\textbf{不合规元标签}并记为 Type~4（流程/字段问题）候选，必须在口径 T 中单独统计、列出比例与示例任务，且不得静默丢弃。若出现整字段缺失（系统错误/字段漏传/导出缺损）则记为 NA，并必须单独统计与报告以供审计。

数学建模：对每个 difficulty tag $k$ 定义二元变量 $z_{t,u,k}\in\left\{0,1\right\}$：若该 tag 被勾选则为 1，否则为 0。\textbf{注意：}在多选界面中，“未被勾选”被操作性定义为 0（该困难未被标注者主张存在），并不要求显式记录每个 tag 的否定项。NA 仅用于整字段缺失（系统错误、字段漏传或提交缺失），必须单独统计与报告以供审计。

Tag 集合（共 6 个）：
\begin{itemize}
    \item 非常简单（trivial）：无明显困难，快速完成
    \item 遮挡明显（occlusion）：家具或人遮挡墙角
    \item 纹理弱（low\_texture）：纯色墙或光照均匀
    \item 拼接缝（seam）：全景拼接缝导致错位
    \item 反光与玻璃（reflection）：镜面伪边界
    \item 画质差或遮罩（low\_quality）：模糊或黑边影响对齐
\end{itemize}

作用：上述 difficulty 标签用于场景分层、路由信号与分布监控。
\subsubsection{Model Issue 多选，仅在半自动条件下收集}
用户界面要求：\texttt{model\_issue} 为多选列表，仅在半自动条件下采集。为避免“空值既可能表示可接受、也可能表示漏填”的不可区分性，界面要求每次提交\textbf{至少选择 1 项}，并施加互斥约束：
\begin{itemize}
    \item 若选择 \texttt{acceptable}（表示初始化质量可接受），则\textbf{不得}同时选择任何其他失效模式；
    \item 若存在任意具体失效模式（如 topology\_failure 等），则\textbf{不得}选择 \texttt{acceptable}。
\end{itemize}
若出现空选、或出现 \texttt{acceptable} 与其他失效模式共存的冲突选择，均视为\textbf{不合规元标签}并记为 Type~4 候选，必须在口径 T 中单独统计、列出比例与示例任务，且不得静默丢弃。若出现整字段缺失（系统错误/字段漏传/导出缺损）则记为 NA，并必须单独统计与报告以供审计。

数学建模：对每个 \texttt{model\_issue} tag $k$ 定义二元变量 $z_{t,u,k}\in\left\{0,1\right\}$：若该 tag 被勾选则为 1，否则为 0。\textbf{注意：}多选界面仅记录被勾选的正例；未被勾选被操作性定义为 0。NA 仅用于整字段缺失（系统错误、字段漏传或提交缺失），必须单独统计与报告以供审计。

Tag 集合（共 8 个）：
\begin{itemize}
    \item 模型标注质量好（acceptable）：无需显著修改
    \item 跨门扩张（overextend\_adjacent）：包含相邻空间
    \item 漏标（underextend）：只标了局部
    \item 过度解析（over\_parsing）：将柱体、墙面凸起、家具边缘等非布局结构误判为真实拐角或边界折点
    \item 角点错位（corner\_drift）：一个或少量角点/局部边界未精确落在真实墙角附近，但拓扑关系仍基本正确
    \item 角点重复（corner\_duplicate）：同一拐角多个点
    \item 拓扑崩溃（topology\_failure）：出现配对失败、闭合失败、自交或其他非法拓扑
    \item 预测失效（fail）：整体严重误导、几乎需从零重画，且无法由上述更具体类型优先概括
\end{itemize}

\subsubsection{改动幅度：$\mathrm{IoU}_{\mathrm{edit}}$}\label{sec:method-iou-edit}
$\mathrm{IoU}_{\mathrm{edit}}$ 用于度量半自动条件下从模型初始化到人工提交的改动幅度（工作量代理），不用于判定正确性。
该指标仅在展示初始化的任务上定义（如 \texttt{PreScreen\_semi}、\texttt{Calibration\_semi} 与 \texttt{SemiAuto\_Test}）；对纯手动条件可记为 NA 并在口径 T 中单独披露缺失率。

\[\mathrm{IoU}_{\mathrm{edit}}\left(t,u\right)=\mathrm{IoU}\left(P_t,A_{t,u}\right)\]
$\mathrm{IoU}_{\mathrm{edit}}$ 越低通常表示改动越大，但不代表改动必然正确，仅用于衡量工作量或改动幅度。
\subsubsection{边界 RMSE}
当角点数量不一致或者配对不稳定时，点对点 RMSE 易失真。为此采用角点生成的上/下边界曲线 $y_{\mathrm{ceil}}\left(x\right)$、$y_{\mathrm{floor}}\left(x\right)$ 在离散网格上的误差作为几何改动幅度指标。
\[\mathrm{BoundaryRMSE}\left(X,Y\right)=\sqrt{\frac{1}{N}\sum_{i=1}^{N}\left[\left(\Delta y_{\mathrm{ceil}}\left(x_i\right)\right)^2+\left(\Delta y_{\mathrm{floor}}\left(x_i\right)\right)^2\right]}\]

\subsubsection{角点匹配 RMSE 辅助指标}
该指标仅在角点配对成功且配对覆盖率足够高的情况下启用。

预测角点的集合为 $\{p_1,p_2,\ldots,p_n\}$，标注角点集合为 $\{g_1,g_2,\ldots,g_m\}$，两点间欧式距离为
\[d\left(p,g\right)=\sqrt{\left(x_p-x_g\right)^2+\left(y_p-y_g\right)^2}\]
用匈牙利算法在两集合间求最小总代价的匹配 $\pi$，匹配对上计算：
\[\mathrm{RMSE}_{\mathrm{match}}\left(t,u\right)=\sqrt{\frac{1}{N}\sum_{i=1}^{N}{d\left(p_i,g_{\pi\left(i\right)}\right)^2}}\]
其中 $N=\min\left(n,m\right)$，取预测角点与标注角点中数量较少的那个。
\subsubsection{活跃标注时间定义}\label{sec:method-active-time}

为明确 RQ1 的主终点度量，特此界定活跃标注时间之操作定义与聚合方式。
\\ 定义与采集条件：
\texttt{active\_time} 指标注者在标注任务页面上的交互活跃时间，而非页面打开挂钟时间。计时器每秒递增 1，当且仅当
\begin{enumerate}
    \item 页面处于可见状态；
    \item 距最近用户交互（鼠标、键盘、滚轮）不超过 \texttt{IDLE\_THRESHOLD} 秒。
\end{enumerate}
当页面不可见超过 \texttt{PAGE\_HIDDEN\_THRESHOLD} 秒应视为中断，需新的交互事件方可恢复计时。阈值在预注册中锁定为 \texttt{IDLE\_THRESHOLD} = 15 s、\texttt{PAGE\_HIDDEN\_THRESHOLD} = 6 s，实验期间不做事后调整。
\\ 聚合协议：
按 session 层级分段记录，其中 session 定义为同一浏览器标签页内的连续交互窗口，以 sessionStorage 的唯一标识追踪。
\begin{itemize}
    \item Per-(task, annotator, session)：取该 session 内的最大累积值，用于消除同一 session 内的重复上报。
    \item Per-(task, annotator)：对跨 session 的多个分段累加，覆盖标注者在不同时段或标签页对同一任务的重复访问。
    \item Per-task：取参与标注的各标注者中的最大值，用于任务级度量。
\end{itemize}

\subsection{多选标签的共识计算}\label{sec:method-tags-consensus}

由于 Difficulty 与 Model Issue 均为多选标签，不存在单一多数类，因此需要专门的共识计算机制。

\subsubsection{基本方案}
采用 Per-tag 二元共识作为主口径；加权共识作为可选干预，见第~\ref{sec:method-weighted-consensus}节，通过预筛选真值锚点与预注册协议规避循环依赖。

共识计算 对每个 tag k 独立：

记号约定：为简洁起见，下述 $c_{t,k}$、$\mathrm{conf}_{t,k}$、$\mathrm{disagree}_{t,k}$ 中任务下标 t 有时省略；实际实现均以任务 t 与 tag k 为索引。

输入：任务 $t$ 的 $m$ 个 worker 标注 $z_{t,1,k},z_{t,2,k},\ldots,z_{t,m,k}$

输出：共识 $c_k$，置信度 $\mathrm{conf}_k$，分歧度 $\mathrm{disagree}_k$

统计勾选人数：$n_{selected}=\sum_{u=1}^{m}\mathbb{I}\left[z_{t,u,k}=1\right]$

统计有效人数：$n_{total}=\sum_{u=1}^{m}\mathbb{I}\left[z_{t,u,k}\in\{0,1\}\right]$（理想情况下$n_{total}=m$；NA 仅用于系统错误、字段漏传或提交缺失，并需单独披露）

若 $n_{total}=0$：$c_k=NA,\ \mathrm{conf}_k=NA,\ \mathrm{disagree}_k=NA$

否则：
$c_k=1$ 若 $\left(n_{selected}/n_{total}\right)\geq0.5$ 否则 0
\[\mathrm{conf}_k=\frac{\max\left(n_{selected},n_{total}-n_{selected}\right)}{n_{total}}\]
\[\mathrm{disagree}_k=1-\mathrm{conf}_k\]

任务级聚合：

共识标签集合 $C_t^{\mathrm{tag}}=\{k:c_k=1\}$

整体分歧度 $\mathrm{disagree}_t=1-\min\{\mathrm{conf}_k: k\in K_t\}$，取最不确定 tag 的不确定性作为保守代理
其中 $K_t=\{k:\mathrm{conf}_k\neq NA,\ n_{total,k}\ge n_{min}\}$，默认 $n_{min}=2$。
若 $K_t$ 为空：令 $\mathrm{disagree}_t=NA$，并单独报告缺失率。

一致性指标 Fleiss' Kappa：
\\ 对每个 tag 独立计算 κ，报告：median κ (across tags)
\\ 期望：κ > 0.5，中等一致性左右（moderate agreement）

若某 tag 的 $\kappa<0.4$：采取预注册处理流程，避免随意切换：
\begin{enumerate}
    \item 频率检查：若该 tag 在 Calibration 中出现频率 < 5\%，按 per-tag 共识 $c_k=1$ 计，则并入其他低频场景，仅用于审计不用于路由分层。
    \item 若频率 $\geq5\%$：不直接用该 tag 作为单独场景做细粒度路由，而是与相近 tag 合并为粗粒度场景，例如 occlusion + low\_texture 合并为低可见性，并在附录中公开合并映射表。
    \item 报告与使用：该 tag 仍在 $Worker\times$ Scene 矩阵与分层表中单独报告，但标注低一致性警告；路由若引用该场景的 $r_u^{\left(s\right)}$，强制使用其置信下界 LCB 而非点估计，避免噪声场景误导。
\end{enumerate}

\subsubsection{可靠度加权共识：核心干预手段}\label{sec:method-weighted-consensus}
动机：结合 Kara 的共识估计与 Liu 的工人分型思想，将可靠与不可靠工人的贡献区分开，以提升标签共识的稳健性与审计可解释性。

适用前提：默认不包含恶意标注者；通过 PreScreen 尽可能滤除随意/不稳定标注者，并在 Discussion 中承认该前提的局限。若 PreScreen 结果显示存在明显恶意或无法稳定区分可靠度，则加权共识不作为主流程决策依据，仅报告未加权多数票与一致性指标用于审计。

关键约束：
权重 $w_u$ 的来源必须独立于当前任务的标签与共识计算。
优先采用预筛选测试样本与专家参考标准得到的初始可靠度 $r_u^{\left(0\right)}$（其构建与专家参考标准见第~\ref{sec:method-prescreen}节），并以此映射到权重 $w_u$；其次采用校准集的 leave-one-task-out 估计。
加权共识用于决策、审计、路由的算法输出，不作为唯一主结论的评价目标；主评价保持独立的参考标准或未加权一致性指标。

加权规则（对每个 tag k 独立）：
令 $z_{t,u,k}\in\{0,1\}$表示 worker u 是否勾选 tag k；权重 $w_u\in\left[0,1\right]$

加权投票率：$p_{t,k}^{\left(w\right)}=\frac{\sum_{u=1}^{m}{w_uz_{t,u,k}}}{\sum_{u=1}^{m}w_u}$

边界条件以避免除零且可审计：若 $\sum_{u=1}^{m}w_u=0$，则该任务在加权共识上记为 NA 并单独报告其发生率；主实现不采用静默 fallback 到未加权多数票，以避免混合 estimand。

加权共识：$c_{t,k}^{\left(w\right)}=\mathbb{I}\left[p_{t,k}^{\left(w\right)}\geq0.5\right]$

权重映射与裁剪：预注册锁定，避免事后调参
\begin{enumerate}
    \item 基础映射：以 PreScreen 测试样本上的初始可靠度 $r_u^{\left(0\right)}$ 为主锚点，映射到 $w_u\in\left[0,1\right]$：
\[w_u=\min\left\{1,\,\max\left\{0,\,\frac{r_u^{(0)}-0.5}{0.5}\right\}\right\}\]
其中$r_u^{\left(0\right)}<0.5$ 视为极端不可靠，设 $w_u=0$。
    \item 上界裁剪：$w_{max}$ 由 Stage 1 PreScreen 的专家参考样本按预注册规则锁定，避免固定常数带来的任意性与事后调参空间。

候选集合：$\mathcal{W}=\{0.80,0.85,0.90,0.95,1.00\}$。
\begin{itemize}
    \item 方案 A（默认，优先）：在 PreScreen 的 20--22 个 manual anchors 上执行重复平衡二折交叉拟合（2-fold repeated balanced split）作为主锁定方案。每次切分均按难度层与主 failure family 保持尽量平衡；在每一折中仅用训练折估计 $r_u^{\left(0\right)}$ 与 $w_u$，再在验证折上评估加权共识与专家参考的一致性增益（$\Delta$ acc 或 $\Delta\kappa$）。3 折分层切分仅作为附录敏感性分析；由于当前 anchor 规模无法稳定支撑每折覆盖，主分析不采用 5 折。

选择规则：取跨折均值最优的 $w_{max}$；若多个候选落在最优的 1-SE 范围内，取其中最小的 $w_{max}$。
    \item 方案 B：若 manual anchors 的覆盖或完成率不足以稳定执行方案 A，则一次性随机拆分 A/B：A 半仅用于估计$r_u^{\left(0\right)}$，B 半仅用于在 $\mathcal{W}$ 上选择 $w_{max}$；拆分随机种子、平衡规则与清单写入预注册并公开。
锁定：最终 $w_{max}$ 在进入 Calibration/Main 前冻结，后续不再调整。
\end{itemize}


    \item 敏感性分析：附录报告 $w_{max}\in\mathcal{W}$ 全网格下结果一致性，用于稳健性披露，不再用于二次选择。
报告方式：同时报告未加权与加权的差异；并在独立参考子集上比较两者与参考标准的一致性与稳定性。
\end{enumerate}

\subsection{质量与可靠度评估：LOO 原则}

可靠度评估遵循两项核心原则：其一，排除自身影响,共识计算不包含被评估者，避免 $r_u$ 系统性偏高；其二，仅用手动标注集估计,避免工具将所有人拉到同一起点。
\subsubsection{任务级一致性 $\mathrm{IAA}$}
对任务 t，收集所有标注者的提交结果$A_{t,u}$ 需要任务标注者数量$\geq2$，定义任务内一致性为两两 IoU 的中位数：
\[\mathrm{IAA}_t=\mathrm{median}_{u\neq v}\,\mathrm{IoU}\left(A_{t,u},A_{t,v}\right)\]
采用中位数而非均值，旨在降低极端标注的影响。$\mathrm{IAA}_t$ 高意味着任务上存在稳定共识，可用于识别专家标注者或建立参考标准。
\subsubsection{共识标注和去除自身后的共识 LOO Consensus}
在任务 t 上，定义共识代表为：
\[C_t=\arg\max_{u}\ \mathrm{median}_{v\neq u}\,\mathrm{IoU}\left(A_{t,u},A_{t,v}\right)\]
为了避免自我参与导致共识虚高，对每个标注者采用去除自身后的共识：
\[C_t^{\left(-u\right)}=\arg\max_{w\in U_t\setminus\{u\}}\ \mathrm{median}_{\substack{v\neq w\\ v\in U_t\setminus\{u\}}}\,\mathrm{IoU}\left(A_{t,w},A_{t,v}\right)\]
并计算：$\mathrm{IoU}_{\mathrm{LOO}}\left(t,u\right)=\mathrm{IoU}\left(A_{t,u},C_t^{\left(-u\right)}\right)$
直接与其他标注者比较：
\[\mathrm{IoU}_{\mathrm{others-med}}\left(t,u\right)=\mathrm{median}_{v\neq u}\,\mathrm{IoU}\left(A_{t,u},A_{t,v}\right)\]

$k=3$ 退化披露：当任务标注者数 $k=3$ 时，LOO 后仅剩 2 人参与共识代表选择，$\mathrm{IoU}_{\mathrm{LOO}}$ 的方差偏大且共识代表退化为二选一。为此，在 $k=3$ 的任务上同时报告 $\mathrm{IoU}_{\mathrm{others-med}}$ 与两两一致性分布作为替代一致性代理；在附录中以该替代代理重复主结论的敏感性分析，预注册锁定为一次，不引入新自由度。
\subsubsection{标注者全局可靠度}\label{sec:method-global-reliability}
标注者 u 在校准集合上的任务集合为 $T_u$
\[r_u=\mathrm{median}_{t\in T_u}\,\mathrm{IoU}\left(A_{t,u},C_t^{\left(-u\right)}\right)\]

估计条件：
\begin{itemize}
    \item 仅在 \texttt{Calibration\_manual} 集合估计
    \item 仅 in-scope 样本
    \item Bootstrap 置信区间：优先采用 BCa bootstrap，重采样次数在预注册中固定为 $B=10,000$；场景分层时采用 stratified bootstrap 保证各场景覆盖。
\end{itemize}
\noindent 说明：本文在实验设计中额外收集 \texttt{Calibration\_semi}（与 \texttt{Calibration\_manual} 同图配对），用于 RQ2 的条件对照（例如 $\mathrm{IAA}_t$ 与反例分布差异）。为避免人的能力估计与初始化质量混杂，$r_u$ 的主估计不使用 \texttt{Calibration\_semi}。
\\ 估计精度停机准则：除最小任务数阈值外，进一步以可靠度估计的置信区间精度作为 Calibration 是否追加任务的客观标准。
\\ 定义：对每个 worker u，计算 $r_u$ 的 95\% CI 半宽
\[h_u=\left(\mathrm{UCB}\left(r_u\right)-\mathrm{LCB}\left(r_u\right)\right)/2\]
其中 $\mathrm{LCB}(r_u)$ 与 $\mathrm{UCB}(r_u)$ 分别表示 $r_u$ 的 95\% 置信区间下界与上界（由本节所述 bootstrap 方案计算）；$\epsilon_r>0$ 为预注册锁定的精度阈值，表示允许的 95\% CI 半宽上限。
\\ 停机规则：若 $h_u\le\epsilon_r$，则该 worker 在 Calibration 阶段不再追加任务；否则仅从预注册的 \texttt{Calibration\_reserve} 任务池中追加小批量任务，并在每次追加后重新计算 CI。\texttt{Calibration\_reserve} 是从预先固定的 \texttt{Calibration\_pool} 中划分出的保留子集，用于避免靠扩大唯一图像池来达标。
\\ 实际执行协议：
\begin{enumerate}
    \item 预先固定 \texttt{Calibration\_pool}（唯一图像集合），并在进入 Calibration 前按预注册随机种子划分为 \texttt{Calibration\_anchor}、\texttt{Calibration\_core} 与 \texttt{Calibration\_reserve}。
    \item 首先发放 \texttt{Calibration\_anchor}，所有通过 PreScreen 的 worker 均完成该集合。
    \item 随后发放 \texttt{Calibration\_core}，采用平衡分配，不要求全员覆盖。收集足够多的多标注数据后，为每个 worker 计算 $r_u$、95\% CI 与 $h_u$。
    \item 对于 $h_u>\epsilon_r$ 的 worker，按预注册批量大小 $\Delta N_{cal}$ 从 \texttt{Calibration\_reserve} 中抽取其尚未完成且可用于 LOO 的任务进行追加，并在每次追加后更新 CI；重复直到满足 $h_u\le\epsilon_r$ 或触达预算上限 $N_{cal,max}$。
    \item 为保证追加任务可用于 LOO 且避免 $k=3$ 退化，追加时优先选择除目标 worker 外已具有至少 3 份其他标注的 reserve 任务（目标 worker 补标后该任务总冗余 $k\ge4$）。若 reserve 内任务尚未满足该条件，则需先按预注册规则补齐其余标注份数，再将该任务用于追加。
\end{enumerate}

\noindent 核心场景覆盖补齐：在完成 \texttt{Calibration\_anchor}+\texttt{Calibration\_core} 的第一轮采集并形成 meta-label 共识后，先按频率与一致性冻结 $\mathcal{S}_{\mathrm{core}}$。对每个 $(u,s)$ 计算缺口
\[\Delta_{u,s}=\max\left(0,\ N_{u,s,\min}-N_{u,s}\right).\]
仅对 $s\in\mathcal{S}_{\mathrm{core}}$ 且 $\Delta_{u,s}>0$ 的对进行覆盖补齐，从 \texttt{Calibration\_reserve} 中选择已形成共识且属于场景 $s$ 的任务补派给缺口较大的 worker。覆盖补齐的目标是启用率与覆盖率达标，不以路由收益为目标。
预算上限：为避免靠无限增加数据才能达标，预注册固定每个 worker 的最大 Calibration 任务数 $N_{cal,max}$（含 reserve）。若达到 $N_{cal,max}$ 仍不达标，则停止追加并标记为 low-precision；后续路由仅使用其 LCB 或直接降级为仅分配低风险任务，并在口径 T 中披露该比例。
\\ 冗余下限：为降低 $k=3$ 下 LOO 的退化与高方差风险，$r_u$ 与 $r_u^{\left(s\right)}$ 的主估计仅使用每任务标注者数 $k\geq4$ 的 \texttt{Calibration\_manual} 任务；$k=3$ 任务仅用于诊断性报告或附录敏感性分析，不作为主估计依据。
\\ 最小任务数阈值：$N_{r,\min}\geq 15$（用于计算 $r_u$ 的最少 \texttt{Calibration\_manual} 任务数；主实现通过 Calibration 任务分配保证该下限可达）。
\\ 独立性补充：初始可靠度锚点 $r_u^{\left(0\right)}$ 由 \texttt{PreScreen\_manual} 的专家参考样本得到（见第~\ref{sec:method-prescreen}节）；后续的 $r_u$ 与 $w_u$ 可视为在此基础上的校准估计，从而打破可靠度与共识的循环依赖。
\\ 说明：为保证 $r_u$ 的估计稳定性与 CI 精度停机准则可执行，主实现通过 \texttt{Calibration\_manual} 的任务分配确保每位通过 PreScreen 的 worker 至少获得约 $N_{cal,\min}$ 个 manual 校准任务（主实现 $N_{cal,\min}=25$）；若因退出或缺失导致不足，则该 worker 标记为 low-precision，并仅以 LCB 或降级路由处理。
\subsubsection{标准 Layout 指标与门控}
为与 HoHoNet、HorizonNet、Dula-Net 等前人研究的常用指标对齐，在满足in-scope 的样本上额外报告以下标准 layout 指标：
\begin{itemize}
    \item 2D IoU 与 3D IoU：HoHoNet 风格的地面投影重合度 2D 与考虑高度体积的重合度 3D
    \item 由布局几何渲染的深度图 RMSE 与 $\delta$：将 layout 按几何关系渲染为全景视角下的每像素深度图，再对两份 layout 的深度图做像素级 RMSE 与阈值准确率
    \item 门控分为两类分别展示：
    \begin{itemize}
        \item 可计算型门控：几何归一化失败、多边形无法构造、layout rendered RMSE 生成失败等
        \item 可比型门控：点对点 RMSE 需要稳定匹配时
    \end{itemize}
\end{itemize}
同时 \texttt{n\_pairs\_mismatch} 不再作为标准 layout 的门控条件，不要求预测和标注角点对数一致，以避免将人工进行显著修正或工作量较大的样本排除。
结构差异相关的信息 $\Delta n_{pairs}$、是否存在奇数个数的点、配对失败等用作分层或诊断的变量，用于报告和反例筛选。

\subsection{结构差异信号}

定义角点对数差：
\[\Delta n_{pairs}\left(t,u\right)=n_{pairs}\left(A_{t,u}\right)-n_{pairs}\left(P_t\right)\]
这是难度与纠错的参考信号，用于分层报告和反例筛选。

结构诊断代理 $g_t$：
为避免未标注前无法获知 difficulty/scene 的情况，本文在任何人工标注之前，仅基于模型初始化/预测结果 $P_t$ 构造结构诊断信号 $g_t$，用于首轮风险分层与触发应激模式。$g_t$ 不刻画正确性，而刻画结构性高风险，典型触发包括：多边形无法构造或自交、拓扑闭合失败、角点重复/异常、以及由 $P_t$ 导出的可计算型门控失败（如几何归一化或渲染失败）。
$g_t$ 的实现以可复现规则为主，阈值与触发集合在进入 Calibration/Main 前预注册锁定；若 $g_t$ 不可计算（如字段缺失），则记为 NA 并在口径 T 中单独报告缺失率。
\subsection{预筛选与工人画像}\label{sec:method-prescreen}

在小样本条件下，少数低质量工人可能显著影响全局结论。因此本文在可靠度估计与路由之前，先通过预筛选识别并剔除不合格或不稳定的工人。
预筛选协议：
\begin{itemize}
    \item 盲测任务池 \texttt{PreScreen\_manual}：从 In-scope 且代表性强的样本中抽取 30 张唯一图像，候选工人完成同一批任务。该任务池不展示模型初始化，其中 20--22 张为带专家参考标准的 manual anchors，用于估计初始可靠度锚点 $r_u^{\left(0\right)}$；其余 8--10 张为随机 non-anchor 样本，用于保留较真实的 base-rate，避免 prescreen 退化为纯 showcase。
    \item 半自动盲测任务池 \texttt{PreScreen\_semi}：从与 \texttt{PreScreen\_manual} 互不重叠的图像集合中抽取约 18 张任务，展示模型初始化作为起点。该任务池由“正常初始化对照子集”与“误导性初始化 trap 子集”构成，用于测量纠错能力与盲信风险，并避免同一图像的重复暴露导致思维惯性。
    \item 锚点样本：\texttt{PreScreen\_manual} 中 20--22 张 manual anchors 用于估计 $r_u^{\left(0\right)}$，但不要求镜像 \texttt{PreScreen\_semi} 的 failure-family 配额结构。
    \item $w_{max}$ 锁定协议：在 \texttt{PreScreen\_manual} 的专家参考样本上按预注册规则锁定并在进入 Calibration/Main 前冻结，具体方案 A/B 与候选集合见第~\ref{sec:method-weighted-consensus}节。
\end{itemize}
\noindent 设计边界：\texttt{PreScreen\_manual} 的目标是测量不依赖初始化的基础标注能力，因此其专家锚点应尽量覆盖主要 difficulty strata，而不要求镜像 \texttt{PreScreen\_semi} 中各类 \texttt{model\_issue} trap 的配额结构。相应的 natural-failure 配额控制仅适用于 semi 盲信子集，而不适用于 manual 锚点池。
\noindent OOS 样本处理边界：若某任务可由专家高一致地判定为 OOS/证据不足，但无法给出稳定的几何参考布局，则其不得进入主几何 GT trap，也不得计入 manual anchors。对此类样本，\texttt{PreScreen\_manual} 仅允许纳入一个固定小配额的 scope-gate 子集，用于检验工人能否识别“不应按正常 in-scope 几何继续标注”的情形；其评分以 scope 判定与拒答/转人工规则是否正确为主，而非几何 IoU。\texttt{PreScreen\_semi} 仅允许纳入更小配额的 OOS stress 子集，用于审计工人面对误导性初始化时是否出现盲信；该子集单独报告，不并入主几何可靠度估计。若专家之间对 OOS 边界本身无法形成高一致判定，则该样本仅进入 ambiguity audit bank，用于附录披露而不作为 hard trap。
测试样本与初始化设计策略：
\begin{enumerate}
    \item \texttt{PreScreen\_manual}：基础能力子集与证据不足子集，用于检验基础操作、稳定性与是否随意猜测。
    \item \texttt{PreScreen\_semi}（纠错子集）：提供约 6 张较高质量的初始化对照样本，检验工人能否在少量修改后得到合理结果。
    \item \texttt{PreScreen\_semi}（盲信子集）：提供约 12 张误导性初始化，优先覆盖跨门扩张、过度解析、局部角点漂移、角点重复等高频失效模式；\texttt{topology\_failure} 与整体 \texttt{fail} 仅保留为小配额稳健性/审计来源，而不作为 prescreen 主 family 配额。
\end{enumerate}
\noindent 误导性初始化的可复现生成：本文以规则化扰动算子为主来源生成误导性初始化。扰动算子集合与强度范围由 Pilot 阶段对常见 \texttt{model\_issue} 的归档抽象得到，并在进入 PreScreen 前冻结。对预先固定的图像子集，先得到模型初始化 $P_t$，再按预注册锁定的算子集合与强度参数将其变换为 $\tilde{P}_t$。每个任务的 \texttt{image\_id}、算子名、强度参数与随机种子均写入清单，并在进入 Main 前冻结与公开，保证可复现与可审计。算子设计以复现常见 \texttt{model\_issue} 为目标，以 \texttt{overextend\_adjacent}、\texttt{corner\_drift}、\texttt{corner\_duplicate} 等在 Pilot 中较多的模式为主；对 \texttt{topology\_failure} 与 \texttt{fail} 等稀有模式仅预注册固定小配额用于稳健性披露。
算子库与冻结字段见附录第~\ref{app:perturbation-library}节。
\begin{itemize}
    \item 输入：固定的图像集合与对应的 $P_t$。
    \item 变换：$\tilde{P}_t=\mathcal{T}_{o,\lambda,\xi}\!\left(P_t\right)$，其中 $o$ 为扰动算子，$\lambda$ 为强度，$\xi$ 为随机种子。
    \item 验收：专家仅对生成结果做验收与记录，剔除明显不可修复或与目标失效模式不一致的样本，专家不参与基于表现的挑选。
\end{itemize}

\noindent 自然失败初始化的补充来源：为提高生态效度，允许在预注册中锁定一小部分 \texttt{PreScreen\_semi} 盲信子集任务来自模型在同分布样本上的自然失败，形成固定的 natural-failure trap bank。该子集与 \texttt{PreScreen\_manual} 图像池严格不重叠，同一 \texttt{base\_task\_id} 不得跨池复用。具体做法是先在固定候选池上运行冻结模型得到 $P_t$，再由专家按锁定的 \texttt{model\_issue} 归档表进行标签化，最后按类型分层随机抽样进入 semi 盲信子集。若某类型自然失败样本不足，则下调该类型 natural-failure 实现配额，并由同类型规则化扰动补足 semi trap 配额；计划配额与实际配额均需披露。该补充来源可用于盲信识别与附录真实性检查，但不得用于设定 manual 锚点、调整扰动算子集合/强度或基于 worker 表现事后挑样，以避免自由度膨胀。
\\ 专家参考标准构建：由 2 名高级标注员独立精标；若两者 $\mathrm{IoU}\geq0.9$ 则取其一致结果作为参考；若 <0.9 则由第 3 人裁决，并将该样本标注为高歧义，用于审计工人在歧义场景下的行为而非作为硬性淘汰依据。
\\ 初始可靠度定义：记 \texttt{PreScreen\_manual} 中带专家参考的任务集合为 $T_u^{(0)}$，其专家参考为 $G_t$，则
\[r_u^{\left(0\right)}=\mathrm{median}_{t\in T_u^{(0)}}\,\mathrm{IoU}\left(A_{t,u},G_t\right)\]
防泄漏原则：PreScreen 与 Calibration 只用校准阶段数据设门槛与估计可靠度，不使用 Validation 或 Test 的任何信息来设阈值或选取样本。
最小门槛为：
\begin{itemize}
    \item 可计算型门控失败率 $\le\tau_{fail}$（可计算性/安全性门槛，不作为 RQ3 主收益主张来源）
    \item Scope 判定一致性 $\geq\tau_{scope}$
    \item 一致性代理 $\geq\tau_r$
    \item 盲信初始化事件比例 $\le\tau_{trust}$
    \item 最小有效任务数 $\geq N_{pre}$
\end{itemize}
\noindent 盲信初始化事件的统计口径：在 \texttt{PreScreen\_semi} 的盲信子集与专家参考子集上，若工人提交与初始化几乎一致且与专家参考明显不一致，则记为一次盲信事件。$\tau_{trust}$ 在 PreScreen 阶段预注册锁定。
\\ 新工人准入：新工人须完成 \texttt{PreScreen\_manual} 与 \texttt{PreScreen\_semi} 并通过门槛，方可进入路由池；该流程写入预注册，以避免事后选择性纳入。
\\ 工人画像 ：
\\ 基本假设与边界：本文默认部署于受控组织环境，主要面对噪声与不稳定标注者。该假设及相应让步见第~\ref{sec:disc-assumptions}节。本文将 spammer 理解为可观测的行为模式，如低质量、随意、系统性偏差或盲信初始化等；若 PreScreen/Calibration 显示异常模式显著，则加权共识与画像在路由中的使用按预注册规则降级，例如仅用于审计或仅用 LCB。

主分析采用“可靠度主轴 + 离散风险轴”的工人画像，而不再把所有过程异常压缩为单一连续分数 $S_u$。该调整的原因是：当前样本规模下，连续异常度对权重与阈值过于敏感，且会混合不同性质的过程量。两轴及其辅助签名均仅使用 PreScreen/Calibration 数据估计，遵循 leave-one-task-out 与时间窗协议，避免循环依赖。
\begin{enumerate}
    \item 可靠度轴（Reliability）：采用全局可靠度的置信下界 $\mathrm{LCB}\left(r_u\right)$ 作为保守排序依据；$r_u$ 的估计、CI 与停机准则见第~\ref{sec:method-global-reliability}节。
    \item 风险轴（Risk tier axis）：定义 worker 的离散风险层级 $R_u^{\mathrm{tier}}\in\{R0,R1,R2,R3\}$，分别对应 Stable、Trust-vulnerable、Condition-fragile、Unusable/Noise。风险层级不由单一加权总分决定，而由三项核心风险指标与两项辅助审计指标共同支撑：
\begin{itemize}
    \item $T_u$：盲信初始化风险。定义来源为 \texttt{PreScreen\_semi}；若工人面对误导性初始化时“几乎不改就提交”，且与专家参考明显不一致，则记为一次 blind-trust 事件。$T_u$ 为该事件在有效 semi 任务中的比例。
    \item $C_u$：条件性脆弱风险。其核心不是全局波动，而是工人在高风险桶上的可靠度显著低于低风险桶或全局水平。形式上以 $\Delta_u=\max_b r_u^{(b)}-\min_b r_u^{(b)}$ 及高风险桶的 $\mathrm{LCB}(r_u^{(b)})$ 共同判断。
    \item $G_u$：可计算失败倾向，即提交导致几何归一化失败、多边形不可构造、渲染失败等工程不可用事件的比例。该量被解释为工程可用性风险，而非恶意风险。
    \item $E_u$：低努力代理，仅作辅助审计，不进入主风险层级的核心判定。定义方式与旧版低努力代理一致，可在图注或 worker card 中披露。
    \item $M_u$：字段缺失/NA 代理，仅作辅助审计，不进入主风险层级的核心判定。定义方式与旧版字段缺失率一致，用于过程质量审计而非主画像轴。
\end{itemize}
上述分量均为过程可观测指标，不依赖主实验的后验性能，因此可用于审计与路由而不构成信息泄露。
\end{enumerate}



功能分组（4 类）：
考虑到本研究的标注者规模上限约 15--20 人，主分析不再使用连续 $S_u$ 与伪精确的四象限式切分，而采用 4 类离散风险层级。这一分组仅依赖 PreScreen/Calibration 的可审计过程指标，并按固定优先级顺序判定。

高风险任务桶 b 的定义仅使用标注前可得代理：$b\in\{g_t\ \mathrm{triggered/not}\}\times\{I_t^{\mathrm{OOD}}=0/1\}$。其中 $d_t$ 与 $I_t^{\mathrm{OOD}}$ 的定义见第~\ref{sec:method-dt}节，$I_t^{\mathrm{OOD}}=1$ 当且仅当 $d_t>\tau_d$，而 $\tau_d$ 由校准集 leave-one-out 分数的预注册分位数锁定。

对每个 b 估计 $r_u^{(b)}$ 及其 LCB。定义与第~\ref{sec:method-global-reliability}节的 $r_u$ 同形式，仅将统计任务集合限制为桶内任务。令 $T_u^{(b)}$ 为 worker u 在桶 b 中的任务集合，则
\[r_u^{(b)}=\mathrm{median}_{t\in T_u^{(b)}}\,\mathrm{IoU}\left(A_{t,u},C_t^{\left(-u\right)}\right)\]
并定义 $\Delta_u=\max_b r_u^{(b)}-\min_b r_u^{(b)}$。为避免小样本下的误判，仅在 $\left|T_u^{(b)}\right|\ge N_{bucket}$ 时使用桶内估计，$N_{bucket}$ 在 PreScreen 阶段预注册锁定。

阈值：$\tau_{r,high},\tau_{r,low},\tau_{trust},\tau_{gap}$ 与高风险桶失效阈值均按分阶段规则锁定：Pilot 锁候选网格与估计协议，PreScreen 锁 admission 与 blind-trust 阈值，Calibration 第一轮后锁依赖场景覆盖的高风险桶启用阈值。

四类风险层级按以下顺序映射：
\begin{itemize}
    \item $R3$（Unusable/Noise）：$\mathrm{LCB}(r_u)\le\tau_{r,low}$，或 $G_u$ 明显偏高，或在 manual anchors 上已无法通过最低 admission。
    \item $R2$（Condition-fragile）：不属于 $R3$，但满足 $\Delta_u\ge\tau_{gap}$，或存在高风险桶使得 $\mathrm{LCB}(r_u^{(b)})\le\tau_{r,low}$。
    \item $R1$（Trust-vulnerable）：不属于 $R3/R2$，且 $T_u\ge\tau_{trust}$。
    \item $R0$（Stable）：不触发上述条件的其余 worker。
\end{itemize}
更细粒度的 6 类画像分型作为未来工作讨论。

\noindent \textbf{顺序判定与可追溯字段。}
为避免多条件重叠导致的歧义，本文采用固定优先级顺序判定：先判定 $R3$；若否，再判定 $R2$；若否，再判定 $R1$；否则归入 $R0$。二维画像图仅用于可视化展示风险层级在“离散横轴 + 连续纵轴”上的分布，不作为分组本体定义。
为保证分组可复现、可审计，在导出与报告侧对每位 worker 记录两个字段：
\begin{itemize}
    \item \texttt{group\_rule\_version}：分组规则版本号（与预注册锁定阈值与触发集合绑定），用于复现实验时精确定位采用的规则；
    \item \texttt{worker\_group\_reason}：导致其进入该组的首要触发原因（与上述顺序判定一致），例如 \texttt{noise\_low\_lcb}、\texttt{noise\_high\_gatefail}、\texttt{fragile\_gap}、\texttt{fragile\_bucket\_fail}、\texttt{trust\_blind}、\texttt{stable\_default}。
\end{itemize}

与路由/共识的接口：
\begin{itemize}
    \item 高风险任务（$I_t^{\mathrm{OOD}}=1$ 或 $g_t$ 触发）：候选池限制为 $R0$，并进一步要求 $\mathrm{LCB}(r_u)\ge\tau_{r,high}$。$R2$ 不进入其已判定失效的高风险桶。
    \item 依赖人工纠错的 semi 任务：默认优先 $R0$，谨慎使用 $R2$，并避免将高 blind-trust 的 $R1$ 作为主力候选池。
    \item 低风险任务：允许 $R0$ 与部分 $R1/R2$ 参与，但需记录其风险签名，并可对高 $E_u$ 或高 $M_u$ 者提高抽检冗余。
    \item $R3$：默认不进入主路由池；若保留则仅用于审计对照，并在共识与加权中按预注册规则置 $w_u=0$。
    \item 过程性证据：在第~\ref{sec:report-process-evidence}节固定抽样的同任务不同功能组对比可视化中，对比不同组 worker 的标注差异，验证分组与路由收益的归因链。
\end{itemize}
\subsection{场景特异可靠度}\label{sec:method-scene-reliability}

工人能力并非全局常数，不同场景类型下同一工人的可靠度可能存在显著差异。因此本文引入场景特异可靠度 $r_u^{\left(s\right)}$，以支持更精细的路由策略。

场景定义：以 difficulty/\texttt{model\_issue} 的共识标签作为场景代理，保留多场景分层作为本文的重要创新点之一（用于异质性审计、路由解释与反例溯源）。为避免小样本下场景过细导致的数据稀疏，本文采用\textbf{预注册的场景选择与合并规则}构造有限场景集合 $\mathcal{S}$：
\begin{itemize}
    \item 候选场景：由 high-level difficulty 与 \texttt{model\_issue} tag 组成。
    \item 不使用 embedding 聚类构造场景；embedding 仅用于 $d_t$ 的 OOD 风险代理与验证集分层。
    \item 两层场景制：主分析仅使用核心场景集合 $\mathcal{S}_{\mathrm{core}}$；其余场景并入 other 作为审计场景集合 $\mathcal{S}_{\mathrm{audit}}$。
    \item 核心场景上限：$\left|\mathcal{S}_{\mathrm{core}}\right|\le S_{\mathrm{core,max}}$，主实现锁定为 $S_{\mathrm{core,max}}=4$。
    \item 核心场景选择规则：仅允许使用频率与一致性，不允许使用任何性能收益或路由收益。具体做法为在完成 \texttt{Calibration\_anchor}+\texttt{Calibration\_core} 的第一轮采集后，统计每个候选场景在 \texttt{Calibration\_manual} 中的出现频率与一致性指标。仅当场景满足 $n_s\ge n_{\min}$ 且 $\kappa_s\ge \tau_\kappa$ 时才进入候选集合，再从中按频率从高到低选择不超过 4 个作为 $\mathcal{S}_{\mathrm{core}}$，其余进入 other。
    \item 罕见但重要场景：如 \texttt{topology\_failure} 与 \texttt{fail}，主分析中仅进入 $\mathcal{S}_{\mathrm{audit}}$，用于频率披露与案例库，不承诺进入 $\mathcal{S}_{\mathrm{core}}$。
\end{itemize}
主实现阈值：$n_{\min}=8$，$\tau_\kappa=0.20$，$\tau_{act}=0.50$，并在进入 Stage 2 前冻结。
估计公式：
\[r_u^{\left(s\right)}=\mathrm{median}_{t\in T_s}\,\mathrm{IoU}\left(A_{t,u},C_t^{\left(-u\right)}\right)\]
其中 $T_s$ 为场景 s 的任务集合。
估计条件：
\begin{itemize}
    \item 只在 \texttt{Calibration\_manual} 估计，保持与全局 $r_u$ 一致。
    \item 使用前置门槛：当 worker 在场景 $s$ 下的有效任务数 $N_{u,s}<N_{u,s,\min}$ 时，不直接使用点估计 $r_u^{(s)}$ 做路由决策。本文保留 $N_{u,s,\min}\in\{5,6\}$ 两档备选，并在 PreScreen 结束后结合观测到的场景频率冻结最终取值。两档门槛下的保守启用率估计表见附录~\ref{app:activation-threshold}。
    \item 置信区间：对满足门槛的 $(u,s)$ 计算 bootstrap 95\% CI，并优先使用其 LCB 作为保守排序依据。
\end{itemize}
\noindent 固定系数收缩（附录敏感性分析，非主实现）：为检验主结论对“跨场景共享先验信息”的鲁棒性，附录额外报告一种不依赖数据驱动调参的固定系数收缩版本。
\[\tilde{r}_u^{\left(s\right)}=\frac{N_{u,s}\cdot r_u^{\left(s\right)}+N_{prior}\cdot r_u}{N_{u,s}+N_{prior}}\]
其中 $N_{prior}$ 在预注册中固定（例如 $N_{prior}=5$），不在实验过程中调参或网格选择。

\noindent 鲁棒性保护（主实现）：若 $N_{u,s}<N_{u,s,\min}$ 或 CI 过宽，则该场景下的路由退化为使用全局 $r_u$（或其 LCB）；若满足门槛，则主实现路由使用 $r_u^{(s)}$ 的 LCB 作为保守排序依据。无论是否退化，均在报告中披露各场景的启用比例与退化比例，避免场景分层在运行中静默失效。收缩版本仅在附录作为敏感性分析报告，且不用于选择或调参。

\noindent 降级声明：若在主报告中观测到总体启用率低于预注册阈值 $\tau_{act}$，则将 $r_u^{(s)}$ 的作用降级为审计与解释性机制，RQ3 的主要收益归因转为全局 $\mathrm{LCB}(r_u)$ 分层与应激模式下的序贯增冗余策略，并继续完整报告 $\mathcal{S}_{\mathrm{audit}}$ 的多场景审计与反例库。

主分析不引入收缩平滑超参数（如 $N_0$），不拟合 $r_u\left(d_t\right)$ 连续函数，也不使用基于 tag 相关性的层级回退。
这些扩展仅放在 讨论板块和未来工作：现有样本规模与周期不足以稳定估计，且会增加调参空间。

\subsection{任务路由策略}

本文不将自演化收敛作为主贡献，而是将路由视为固定预算下的质量控制策略：在固定预算内，将更高风险的任务优先交由更稳定的工人，并用序贯追加与停止准则把冗余花在刀刃上。核心要求是可审计、可离线模拟，且在 non-IID 与 OOD 条件下保持稳健。

为打破冷启动困境，本文将路由信号区分为标注前可得与标注后可得两类，并将二者的使用时点写入预注册：
\begin{enumerate}
    \item 标注前可得的风险代理（cold-start proxies）：
    \begin{itemize}
        \item 分布风险 $d_t$：由 HoHoNet 共享 pre-head latent 经宽度方向全局平均池化与 L2 归一化后，相对于校准参考库的 kNN 距离得到；连续分数 $d_t$ 用于审计与排序，二值指示 $I_t^{\mathrm{OOD}}$ 用作触发器与分层变量。
        \item 结构诊断 $g_t$：由模型初始化/预测结果 $P_t$ 的结构合法性与失败模式构造，用于识别结构性高风险。
    \end{itemize}
    \item 标注后可得的人的诊断信号（post-submission diagnostics）：
    $C_t^{\mathrm{tag}}$：difficulty 与 \texttt{model\_issue} 的共识标签集合；以及基于 per-tag $\mathrm{conf}_k$ 的任务级分歧度 $\mathrm{disagree}_t$。
    这些信号仅在任务已获得至少 2 份有效元标签后启用，用于序贯追加与解释性分层，不作为首轮派单的前置条件。
    \item 系统失败信号：门控失败类型、字段缺失率（NA 缺失单独审计）。
\end{enumerate}

预注册主实现（$d_t$）：\label{sec:method-dt}
embedding 来源：固定采用 HoHoNet 共享 pre-head 1D latent $h_t$（即 layout head 之前的共享水平特征），不在本文中微调；将其沿水平宽度 $W$ 做全局平均池化得到图像级向量 $z_t=\mathrm{GAP}_{W}(h_t)$，再做 L2 归一化 $\tilde z_t=z_t/(\|z_t\|_2+\varepsilon)$。若因环境原因主特征不可用，仅在附录给出替代 embedding 并报告其与主实现 $d_t$ 的一致性。
\\ 参考库：固定用 \texttt{Calibration\_manual} 构造校准参考库 $\mathcal{B}_{\mathrm{cal}}=\{\tilde z_i\}_{i\in \mathcal{C}}$，不混入 validation/test；校准阶段样本自身的分数一律采用 leave-one-out 方式计算，避免自邻近导致阈值偏乐观。
\\ $d_t$ 定义：主实现遵循 KNN-OOD 的主分数形式，采用归一化特征到校准参考库的第 $K$ 个近邻欧氏距离
\[d_t = D_K\!\left(\tilde z_t,\mathcal{B}_{\mathrm{cal}}\right)=\left\|\tilde z_t-\tilde z_{(K)}\right\|_2\]
其中 $\tilde z_{(K)}$ 表示 $\tilde z_t$ 在参考库中的第 $K$ 个最近邻。主实现固定 $K=10$；附录做 $K\in\{5,10,20\}$ 的敏感性分析。
\\ 阈值与触发器：记校准集样本的 leave-one-out 分数为 $\{d_i^{\mathrm{cal,LOO}}\}$，主实现固定 $q=90\%$，并定义
\[\tau_d = Q_q\!\left(\{d_i^{\mathrm{cal,LOO}}\}\right),\qquad I_t^{\mathrm{OOD}}=\mathbf{1}[d_t>\tau_d]\]
其中 $d_t$ 为连续偏移风险分数，$I_t^{\mathrm{OOD}}$ 为预注册的高风险触发器。附录额外报告分位数 $q\in\{85,90,95\}$ 与“前 $K$ 个近邻平均距离”替代打分的敏感性分析，但二者均不作为主实现。
\\ 敏感性分析：Mahalanobis 距离仅作为附录替代度量，并在与主实现相同的 pooled、归一化特征上计算；主分析不引入 Lee 等人原始 OOD detector 中额外的逻辑回归检测器，以避免超出本文的主要 estimand。
\\ 严谨定位：$d_t$ 与 $I_t^{\mathrm{OOD}}$ 不直接判对错，不作为硬门控过滤样本；它们仅用于审计偏移是否存在、触发应激模式、以及在第~\ref{sec:exp-iid-stress}节中做分层对比。

策略设计遵循以下原则：
\begin{itemize}
    \item 预筛选约束：仅在通过 PreScreen 的工人池中路由；未通过者只分配低风险任务或剔除。
    \item 保守选择：用 $r_u^{\left(s\right)}$ 的置信下界进行排序，避免小样本场景下点估计误导。
    \item 反 domination：引入负载惩罚或连续派单上限，避免少数工人被过度使用。
    \item 应激模式（stress mode）：当 $I_t^{\mathrm{OOD}}=1$、$g_t$ 触发或系统失败信号异常时进入应激模式，收紧候选工人至高置信下界子集，并提高冗余上限或触发专家复核队列。
\end{itemize}
序贯冗余与停止准则：
\begin{itemize}
    \item 首轮派单：按风险分层设置起始冗余 $k_0\left(t\right)$。低风险任务默认 $k_0=1$。高风险任务默认 $k_0=2$，高风险由 $I_t^{\mathrm{OOD}}=1$ 或 $g_t$ 触发定义。应激模式下将 $k_0$ 提升为 3，并收紧候选池至更高 $\mathrm{LCB}(r_u)$ 的子集。
    \item 追加触发：当任务累计获得至少 2 份有效元标签后，计算 $\mathrm{disagree}_t$。若 $\mathrm{disagree}_t$ 高于阈值、出现门控失败或字段缺失异常、或有效人数不足，则每次追加 1 份标注。
    \item 停止条件：当 $\mathrm{disagree}_t$ 低于阈值，或等价地最不确定 tag 的 $\mathrm{conf}_k$ 达标，且字段缺失不触发审计异常时停止追加。
    \item 硬上限与降级：预注册固定最大冗余 $k_{\max}$，主实现锁定为 $k_{\max}=5$。若达到 $k_{\max}$ 仍未达标，则不再追加，进入专家复核与单独审计队列，并在口径 T 中披露该类任务比例。
\end{itemize}
本文对比以下三种策略：
\begin{itemize}
    \item 策略 1：Random baseline，不使用画像、场景、OOD。
    \item 策略 2：Global routing，仅用全局 $r_u$。
    \item 策略 3：PreScreen + Stratified + OOD-aware routing，路线 A 主体。
\end{itemize}

离线模拟评估协议：
\begin{itemize}
    \item 数据：三策略主对比使用 \texttt{Calibration\_manual} 的多标注任务（每图 $k=4$；含 \texttt{Calibration\_anchor} 全员覆盖），以保证在不做反事实插补的前提下可回放策略 1/2/3 在同任务、同候选池下的选择与序贯停止。$Validation_{semi}$ 若部署仅执行策略 3，则用于报告策略 3 的实际运行表现与 non-IID/OOD 应激收益；策略 1/2 在该集合上不作为主对比。
    \item 方法：模拟预筛选→序贯分配→停止，并取真实提交作为输出，不做事后挑选。
    \item 影子策略（可选补充披露）：若在 $Validation_{semi}$ 上额外报告策略 1/2 的“离线回放”结果，仅在其决策落入已观测标注者集合的支持集上计算指标，并同时披露支持率；不可支持部分不得以任何模型插补补齐。
    \item 防后见之明协议：
    \begin{itemize}
        \item Leave-one-task-out：对每个任务 t，估计 $w_u$ 与 $r_u^{\left(s\right)}$ 时不使用 t 的任何信息。
        \item 时间顺序：在数据采集阶段记录每次提交的时间戳；离线模拟时按时间构造历史窗口，仅用 t 之前的历史任务估计 $w_u$ 与 $r_u^{\left(s\right)}$ 再对 t 进行路由。
        \item 时间戳缺失处理：若时间戳字段缺失，记录为 NA 并在口径 T 中单独报告；该任务不进入时间窗口敏感性分析，但仍可进入 LOO 协议下的主分析。
        \item 冷启动：当某 worker 的历史有效任务数 <$N_{hist,\min}=5$ 时，其 $r_u$ 与 $r_u^{\left(s\right)}$ 不启用，路由退化为在通过 PreScreen 的池内随机或仅按 $r_u^{\left(0\right)}$ 粗排；当历史 $\geq$ $N_{hist,\min}$ 后再启用基于 $r_u$ 与 $r_u^{\left(s\right)}$ 的分层路由。$N_{hist,\min}$ 存在离散跳变效应，即阈值附近的标注者路由策略会发生突变；为证明结论不依赖该特定阈值，预注册固定比较 $N_{hist,\min}\in\{3,5,7\}$ 三档的敏感性分析（仅影响路由启用门槛，不改变 $r_u$ 的估计口径）。
    \end{itemize}
    \item 指标（主）：一致性代理、预算效率 $\bar{k}_{\mathrm{used}}$（或等价预算-覆盖关系）。
    \item 安全性/数据可用性审计（非主终点）：门控失败率与失败原因分布、关键字段缺失率；用于证明不静默过滤与口径 T/I/M 的可追溯性，不作为“路由收益”的主证据。
\end{itemize}
\subsection{OOS（Out-of-Scope）统计与混合判定}\label{sec:method-oos}

对于标记为 OOS 的样本，主指标分析中予以剔除，但单独报告其比例、子类分布与判定一致性。
混合判定：
\begin{itemize}
    \item 当存在 u 选 In-scope 且 v 选 OOS 时，标记为 \texttt{scope\_conflict}
    \item 人工复核协议：由两名独立高级标注员盲法复核；若分歧则由第三人裁决，并报告需要第三方裁决的比例。
    \item 报告混合判定比例
\end{itemize}

\subsection{反例筛选规则}\label{sec:method-counterexample}

本文定义三类反例，各类均有明确的可复现条件：
\begin{enumerate}
    \item \textbf{Type 1：低一致性}
    \begin{itemize}
        \item 条件：$\mathrm{IAA}_t<0.7$ 且 $n_{annotations}\geq3$
        \item 可能原因：任务歧义或标注错误
    \end{itemize}
    \item \textbf{Type 2：异常编辑模式}
    \begin{itemize}
        \item 2a 改动极大但 $\mathrm{IAA}$ 未提升：$\mathrm{IoU}_{\mathrm{edit}}<0.5$ 且 $\mathrm{IAA}_t<\mathrm{median}(\mathrm{IAA})$
        \item 2b 改动极小但与他人严重不一致：$\mathrm{IoU}_{\mathrm{edit}}>0.95$ 且 $\mathrm{IoU}_{\mathrm{LOO}}<0.7$
    \end{itemize}
    \item \textbf{Type 3：几何合法性失败}
    \begin{itemize}
        \item 条件：提交结果出现自交、多边形无法构造、拓扑闭合失败，或其他预注册的几何合法性违规
        \item 说明：该类反例针对结构非法本身，而不是一般性的 metric gate/导出缺失；后者另作安全性与数据可用性审计披露
    \end{itemize}
\end{enumerate}
人工复核流程分为四个步骤：
\begin{enumerate}
    \item 自动筛选符合规则的候选；
    \item 专家逐一复核；
    \item 确认反例类型；
    \item 导出案例库并附带可视化。
\end{enumerate}



\section{实验设置}

实验旨在同时验证三个研究问题：半自动初始化的效率收益（RQ1）、对质量与可靠性分布的影响（RQ2）、以及场景感知路由的收益（RQ3）。以下依次介绍数据组织、评估协议与统计检验方案。
\subsection{IID 与 non-IID 应激测试}\label{sec:exp-iid-stress}

为检验审计与路由机制在分布偏移下的价值，本文在同一套标注预算与评估协议下同时覆盖两类部署情景。
IID baseline：训练、校准、验证来自近似同分布的数据抽样，采用常规随机抽样与分层控制。

non-IID stress test：通过数据切分显式制造难度与失效模式的分布偏移，用于检验审计与路由机制在偏移下是否更有价值。

构造方式：
为避免用标注后才产生的元标签反向决定数据切分，本文以标注前可得的代理信号构造应激情景：在验证集中提高高风险代理任务的占比，其中高风险由 $I_t^{\mathrm{OOD}}=1$ 与 $g_t$ 的触发共同定义。这里 $I_t^{\mathrm{OOD}}=\mathbf{1}[d_t>\tau_d]$，且 $\tau_d$ 仅由校准集 leave-one-out 分数的预注册分位数估计，不使用验证集结果调参。

同时，为保证 non-IID 的语义可解释性，在标注后再用 difficulty/\texttt{model\_issue} 的共识标签作审计性验证：定义困难高风险集合 $S_{hard}$（如 occlusion、low\_texture、corner\_duplicate、corner\_drift、overextend\_adjacent 等 tag），并报告其在 IID 与 non-IID 切分中的占比差异；该验证用于解释性披露，而非用于切分本身。

\texttt{Validation\_OOD} 定义：可复现、避免硬门控，在验证集中取满足 $I_t^{\mathrm{OOD}}=1$ 的子集，等价于 $d_t>\tau_d$；主实现锁定为校准集 leave-one-out 分数的 $q=90\%$ 分位阈值。该阈值在 IID 近似条件下对应约 10\% 的预期高风险比例，但在真实 validation 或 non-IID 条件下其实际占比允许偏离 10\%。附录做 $q\in\{85,90,95\}$、$K\in\{5,10,20\}$ 与“前 $K$ 个近邻平均距离”替代打分的敏感性分析。

Hard 子集 H 与 \texttt{Validation\_OOD} 采用不同的构造口径：\texttt{Validation\_OOD} 仅由 $I_t^{\mathrm{OOD}}=1$ 定义；H 允许由 $g_t$ 触发纳入，因此 H 与 \texttt{Validation\_OOD} 通常为部分重叠关系。本文在结果中分别报告二者的规模与重叠比例，并在对应子集上重复主要对比，以避免将两者混作同一集合。

IID：保证训练与验证中由 $I_t^{\mathrm{OOD}}$/$g_t$ 定义的高风险代理占比接近，或在预先声明的误差范围内；并报告事后审计得到的 $S_{hard}$ 占比作为一致性检查。

non-IID：在验证集中人为提高高风险代理任务（$I_t^{\mathrm{OOD}}=1$ 或 $g_t$ 触发）的比例，使其在该验证集内占比显著偏高；训练集与校准集则保持较低的高风险代理占比。随后用标注后的 $S_{hard}$ 占比与失败模式分布作审计性验证，说明该切分确实带来更高的困难/失效暴露。
核心比较：
\begin{itemize}
    \item OOD-aware 审计任务：报告 $d_t$ 的分布、$I_t^{\mathrm{OOD}}$ 的触发比例，并给出 $d_t$ 与门控失败、反例率、$\mathrm{IAA}_t$ 的相关性（相关而非因果），用于说明分布偏移的存在，以及偏移条件下系统更容易失效。
    \item 防选择偏差：无论 $d_t$ 高低，样本仍进入口径 T/I/M 的透明披露；$I_t^{\mathrm{OOD}}$ 仅用于分层报告与策略触发，不用于静默过滤。
\end{itemize}
\subsection{条件分组与数据组织（$\mathbf{N}_{\mathrm{pool}}=458$）}\label{sec:exp-data-org}

记号：$N_{img}$ 为唯一图像数，$N_{task}$ 为任务实例数，k 为每图标注份数，$\left|U\right|$ 为 PreScreen 候选工人数。
\noindent 序贯派单口径：$k_0$ 为起始冗余，$k_{\max}$ 为最大冗余，$k_{\mathrm{used}}$ 为每任务实际使用标注份数。

主方案分组：
\begin{itemize}
    \item Stage 0 Pilot：$N_{img}=15$，$k=2$；用于流程验证、\texttt{model\_issue} 归档与协议锁定，不进入主结论。Pilot 阶段锁定的不是最终 admission 数值，而是不可变协议边界：manual/semi/OOS 三池边界、同一 \texttt{base\_task\_id} 不跨池复用、manual anchor 选择规则、semi trap schema、扰动算子集合与强度档位、候选权重网格与 split 协议、审计字段与随机种子策略。
    \item Stage 1 PreScreen：\texttt{PreScreen\_manual} 与 \texttt{PreScreen\_semi} 使用两套互不重叠的图像池。\texttt{PreScreen\_manual} 固定为 $N_{img}=30$，其中 20--22 张为带专家参考标准的 manual anchors，用于估计 $r_u^{\left(0\right)}$、锁定 admission 与 $w_{\max}$；其余 8--10 张为随机 non-anchor 样本，用于保留 base-rate 并避免全由 showcase 题构成。对专家可高一致判定为 OOS/证据不足但无法给出稳定几何参考的样本，仅允许以固定小配额进入 manual 的 scope-gate 子集，按 scope 判定正确性计分，而不作为几何 GT anchor。\texttt{PreScreen\_semi} 固定为约 $N_{img}=18$，用于测量纠错能力与盲信风险；其结构为“正常初始化对照子集 + 误导性初始化 trap 子集”，误导性初始化采用双来源但主次分明：主来源为规则化扰动算子并冻结清单与随机种子，补充来源为固定小配额的自然失败 trap bank，仅进入 semi 盲信子集；此外仅允许更小配额的可判定 OOS stress 子集用于盲信审计，并与主几何 trap 分开报告。若某类自然失败样本不足，则下调其实际配额，并由同类型规则化扰动补足 semi trap 配额；计划配额与实际配额一并披露。对专家之间本身高歧义的 OOS 样本，仅进入附录审计库，不作为 hard trap。
    \item Stage 2 Calibration：\texttt{Calibration\_manual} 固定校准池 $N_{img}=100$，每图冗余 $k=4$，用于估计 $r_u$ 与 $r_u^{\left(s\right)}$（见第~\ref{sec:method-global-reliability}节与第~\ref{sec:method-scene-reliability}节）。该阶段不要求每位 worker 标注每张图，而是采用平衡任务分配使每位通过 PreScreen 的 worker 至少获得 $N_{cal,\min}=25$ 个 manual 校准任务，并与 CI 精度停机准则结合按需从 \texttt{Calibration\_reserve} 追加。
    \\
    为提高 worker 间可比性并稳定 LOO 共识，\texttt{Calibration\_manual} 在唯一图像池 100 内进一步划分为三块：
    \begin{itemize}
        \item \texttt{Calibration\_anchor}：$N_{img}=12$，所有通过 PreScreen 的 worker 均完成。
        \item \texttt{Calibration\_core}：$N_{img}=75$，每图 $k=4$，采用平衡分配，不要求全员覆盖。
        \item \texttt{Calibration\_reserve}：$N_{img}=13$，仅用于追加任务。
    \end{itemize}
    \texttt{Calibration\_reserve} 的用途预注册写死为两类：其一为 CI 半宽未达标时追加；其二为核心场景启用率不足时的覆盖补齐。不得以提升路由收益为目标使用 reserve。
    \\
    \texttt{Calibration\_semi} 仅在从上述 100 张中预注册抽取的配对子集 $N_{img}=25$ 上收集，且每图冗余 $k=4$，用于 RQ2 的条件对照。为避免混杂，\texttt{Calibration\_semi} 不用于估计 $r_u$ 或 $r_u^{\left(s\right)}$。manual/semi 的顺序随机化。
    \item Stage 3 Main：
    \begin{itemize}
        \item RQ1 Test：\texttt{Manual\_Test} 与 \texttt{SemiAuto\_Test} 同图 $N_{img}=100$，$k=1$。
        \item RQ3 Validation：$Validation_{semi}$，$N_{img}=60$，起始冗余 $k_0=3$，最大冗余 $k_{\max}=5$，采用序贯追加与停止；$Gold_{manual}$，$N_{img}=20$，$k=3$ 用于稳健性对照。
    \end{itemize}
\end{itemize}
预注册约束：
\begin{itemize}
    \item 带专家参考标准的样本与 Main 物理隔离；Calibration/Validation/Test 图像不重叠；Pilot/PreScreen 不进入 Main 对比。
    \item Hard 子集 H 从 Validation 构造，要求$\left|H\right|\geq30$；入选依据为 $I_t^{\mathrm{OOD}}=1$ 或高 $g_t$，并在标注后验证 $S_{hard}$ 占比$\geq70$。
    \item OOD 仅用于分层与触发，不用于静默过滤；全部样本进入 T/I/M 披露。
    \item 规模说明：核心主分析唯一图像约 260（Calibration\_manual 100 + Validation 60 + Test 100），其余用于前置锚点、流程验证与稳健性。
\end{itemize}
数据与评测口径锁定：
\begin{itemize}
    \item test 与 \texttt{test\_no\_occ} 的 img 集合一致，差异仅在 \texttt{label\_cor} 修正。
    \item 标准 benchmark 主报告使用 test/\texttt{label\_cor}；\texttt{test\_no\_occ}/\texttt{label\_cor} 作为附录敏感性分析。
\end{itemize}
\subsection{参与者与 worker mix 控制}

培训内容包括 Label Studio 基础操作、仅标注相机所在房间的规则、以及 OOS 分流规则。
\\ 分配与记录：统一培训后不再引入新手与熟手分层。为控制 worker mix 与能力差异混杂，采用 worker 级别的平衡分配：每位通过 PreScreen 的 worker 在 Manual 与 Semi 两种条件下分配到近似相同数量的任务，并随机化任务顺序。同时记录每个 worker 的任务数、完成时间与返工率，并报告每批 worker 组成以供 worker mix 控制。
新工人准入按第~\ref{sec:method-prescreen}节执行。

\subsection{元标签采集协议与提交时硬校验（Type 4 预防）}

本文将元标签（\texttt{scope}/difficulty/\texttt{model\_issue}）定位为实验协议的一部分，而非事后清洗时的自由度。目的在于避免 ``空值=无困难'' 与 ``漏填'' 不可区分，从而提升采集系统的可控性与可审计性。

我们采用\textbf{两层防线}：
\begin{itemize}
    \item \textbf{前端提交时硬阻断（UI-level hard validation）}：Label Studio 视图配置中对关键字段启用 \texttt{required="true"}（至少选 1 项），从而使 ``空选'' 在交互层面无法提交；同时对多选字段加入互斥规则（difficulty 的 \texttt{trivial} 与其他困难标签；\texttt{model\_issue} 的 \texttt{acceptable} 与其他失败模式），在用户点击 Submit/Update 时由前端脚本即时阻断提交并提示冲突。
    \item \textbf{导出/落盘前兜底校验（backend guard for audit）}：对导出的 Label Studio JSON 再次执行同一组规则校验，将不合规记录写入拒收清单并从主分析输入中剔除，同时保留拒收原因分布作为审计披露。这一层用于覆盖 ``前端不可控''（版本差异、脚本未加载、导出缺损等）导致的残余风险。
\end{itemize}

现实边界：上述机制可将 ``空选/互斥冲突'' 的发生率压到接近 0，但无法保证 Type 4 绝对为 0——系统侧仍可能出现 NA（字段未传输、导出损坏、版本不一致）。因此本文仍将 NA 与残余不合规事件作为口径 T 的审计项披露，并在结果中给出明确分母与示例任务，而非将其 ``从论文中消失''。
\subsection{预注册协议}

\subsubsection{RQ1 主终点：效率}
\begin{itemize}
    \item 指标：\texttt{active\_time} 中位数差异（Manual vs Semi）
    \item 统计：Mann--Whitney $U$ 检验（非参数）
    \item 效应量：Hodges–Lehmann 中位数差
    \item CI：95\% CI，优先 BCa bootstrap，重采样次数预注册固定
    \item worker mix 控制：通过前文参与者与 worker mix 控制中的 worker 级别平衡分配，减少条件间能力差异带来的混杂；结果中报告两条件下 worker 构成与任务覆盖的平衡性。
    \item 风险控制：回归分析 $time \sim \mathrm{IoU}_\mathrm{edit}$，检验时间缩短并非源于随意修改
    \item 一致性校验：开展自定义计时与原生计时的一致性校验，计算自定义统计的 \texttt{active\_time} 与 Label Studio 原生 \texttt{lead\_time} 的 Spearman 相关系数，同时分析二者比值的中位数及四分位距（IQR），以此验证自定义活跃标注时间的采集逻辑无系统性偏离。该校验成本较低，可排除计时采集环节的系统性实现误差。
\end{itemize}

\subsubsection{RQ2 主终点：可靠性}
\begin{itemize}
    \item 指标：在预注册的同图配对子集（$N_{img}=25$）上比较 $\mathrm{IAA}_t$ 的差异：\texttt{Calibration\_manual} vs \texttt{Calibration\_semi}
    \item 统计：同图配对置换检验。对每张配对图像定义 $\Delta_t=\mathrm{IAA}_t^{semi}-\mathrm{IAA}_t^{manual}$，在零假设下随机翻转 $\Delta_t$ 的符号，得到 $\Delta_t$ 的和或中位数差的置换分布。
    \item 效应量：$\mathrm{median}(\Delta_t)$，并用配对 bootstrap 给出 95\% CI。
    \item 反例数量：在同一配对子集上比较三类反例分布差异（$\chi^2$ 检验）
    \item RQ2b：加权共识增益，口径见第~\ref{sec:method-weighted-consensus}节。
    \item 评估集合：\texttt{PreScreen\_manual} 的专家参考子集及其扩展复核子集。$w_{max}$ 的锁定与性能评估采用预注册的方案 A，样本不足时使用降级方案 B；在每个验证折上，$r_u^{(0)}$ 与 $w_u$ 的估计均不使用该验证折的标签，保证评估独立性。
    \item 主终点：加权与未加权共识对专家参考标准的一致性提升，如 per-tag accuracy 或 Cohen's $\kappa$ 的提升。
    \item 统计：按任务重采样的 BCa bootstrap 得到 $\Delta\kappa$ 或 $\Delta\,\mathrm{acc}$ 的 95\% CI；若 CI 不跨 0 且效应量 $\geq MDE_\kappa$（默认 0.05），则判为工程上有意义的提升。
    \item 反例检出：对第~\ref{sec:method-counterexample}节的反例规则产生的候选集，比较加权与未加权的检出率差异，使用 McNemar's test 或 bootstrap CI。
    \item 补充报告：在附录中补充 Kolmogorov--Smirnov 检验作为分布形状差异的敏感性分析，不作为主终点。
\end{itemize}
\subsubsection{RQ3 主终点：路由}
\noindent 数据来源与可比性说明（解决只实际跑策略 3 时，策略 1/2 从哪来）：
\begin{itemize}
    \item \textbf{主对比数据（离线回放 / shadow policy evaluation）}：三策略的可比性主依赖 \texttt{Calibration\_manual} 的多标注任务（\texttt{Calibration\_core} 每图 $k=4$，并含 \texttt{Calibration\_anchor} 全员覆盖）。该阶段的任务分配为预注册的平衡分配，与策略 1/2/3 无关，因此可将每图实际获得的标注者集合视为该任务的候选池 $U_t$，并在 \emph{不做反事实插补} 的前提下回放三种策略在 $U_t$ 内的选择与序贯停止，从而得到同任务、同候选池下的可审计对比。
    \item \textbf{部署数据（单一 Validation）}：$Validation_{semi}$ 用于报告策略 3 的实际运行表现与 non-IID/OOD 应激收益（安全性事件比例：门控失败/字段缺失等；反例拦截率；预算效率等）。若部署时仅执行策略 3，则策略 1/2 在该集合上不作为主对比；如需补充影子策略结果，仅在策略 1/2 的决策落在已观测标注者集合的支持集上报告，并同时披露支持率（不可支持部分不得以任何模型插补补齐）。
\end{itemize}
\begin{itemize}
    \item 指标 1 主：一致性代理提升（离线回放）：对每任务在候选池 $U_t$ 内按策略选择标注者 $u$，计算 $\mathrm{IoU}\left(A_{t,u},\ C_t^{(-u)}\right)$（以同任务其余标注构造的 LOO 共识 $C_t^{(-u)}$ 为参考）。
    \item 统计 1：同任务配对比较（优先配对置换或配对 bootstrap CI；t-test 仅作为附录敏感性分析），仅在 $\left|U_t\right|\geq3$ 的任务上计算。
    \item 效应量：Cohen's d 标准化的均值差效应量
    \item 指标 1b：预算效率，序贯停止，平均每任务使用标注份数 $\bar{k}_{\mathrm{used}}$ 或在固定预算下可覆盖的任务数
    \item 安全性/数据可用性审计（非主终点）：门控失败率与失败原因分布、关键字段缺失率；用于证明不静默过滤与口径 T/I/M 的可追溯性。原则上不将“失败率降低”作为路由收益主证据，也不将其纳入主检验集合。
    \item 应激测试报告（non-IID）：在前文定义的 \texttt{Validation\_OOD} 子集上重复上述对比；主证据仍以一致性代理与预算效率为核心，并单独披露安全性事件分布（门控失败/字段缺失等）作为风险信号。
\end{itemize}
\subsection{多重检验处理}

策略：预先限制检验集合，避免 FDR 复杂度。
\begin{itemize}
    \item RQ1：1 个主检验 Manual vs Semi = 1 个
    \item RQ2：1 个主检验（$\mathrm{IAA}_t$ 配对置换） + 1 个反例分布检验（$\chi^2$） + 2 个 RQ2b 检验（加权共识增益、反例检出差异） = 4 个
    \item RQ3：2 个主检验 一致性代理 / 预算效率 = 2 个
    \item 总计：7 个计划检验
    \item 校正：Bonferroni $\alpha = 0.05/7 \approx 0.0071$（保守）
\end{itemize}
\subsection{效应量门槛与 MDE}

最小可解释差异（MDE，示例阈值，以试运行与先验为准）：
\begin{itemize}
    \item Time：$\geq8\%$ 相对改进
    \item 一致性代理：$\geq0.05$
    \item 预算效率：预注册的最小可解释改进（例如平均 $\bar{k}_{\mathrm{used}}$ 的相对减少或等价覆盖提升）
\end{itemize}

报告原则：只有同时满足 $p<0.05/7$（计划检验 Bonferroni 校正后）且效应量超过 MDE 时，才声称显著改进。

MDE 的确定采用分级锁定流程：
\begin{enumerate}
    \item Stage 0 Pilot：预研，$N_{img}=15$，估计工具可用性与噪声底线，并锁定协议骨架而非最终 admission 数值：manual/semi/OOS 三池边界、anchor/trap 抽样框架、\texttt{PreScreen\_semi} 的算子集合与强度范围、候选阈值网格、A/B 或 2-fold split 协议、审计字段与随机种子策略。
    \item Stage 1 PreScreen：在独立锚点与 non-anchor 样本上估计工人表现方差与可达上限，并将候选 $\tau_0$ 调整并锁定为 admission 相关阈值与最终 $w_{max}$。主锁定对象包括：通过工人名单、$w_{max}$、$\tau_{fail}$、$\tau_{scope}$、$\tau_r$、$\tau_{trust}$，以及是否启用加权共识主实现或降级为未加权审计。
    \item Stage 2 Main Experiment：Calibration 第一轮后，才最终冻结依赖场景覆盖与频率的 scene-routing 参数，例如核心场景集合 $\mathcal{S}_{core}$、$N_{u,s,\min}$、场景启用/退化规则与高风险桶阈值；PreScreen 与测试样本不得进入 RQ3 的路由效果评估集合，防数据泄漏。
\end{enumerate}

实施说明：三段式流程允许软件未完全自动化；可通过人工分发不同任务包完成阶段切换，但必须保证不手动修改标注内容，且阶段边界与样本隔离可审计。



\section{报告与可审计输出}

本文的报告体系以可审计性为核心，通过三级口径、分层矩阵与反例库构成完整的证据链。
\subsection{三级口径 T/I/M}

为避免因指标不可计算而静默丢弃样本，本文采用三级口径透明披露所有数据：
\begin{itemize}
    \item 口径 T（Total）——透明披露：包含所有收集到的标注，不因指标无法计算而静默丢弃。报告 OOS 比例、门控失败原因分布、缺失率。
    \item 口径 I（In-scope）：仅包含 \texttt{scope}=\texttt{in-scope} 且 \texttt{scope} 非缺失的样本，用于主指标分析（效率、IAA、$r_{u}$）。OOS 样本单独报告比例与子类。
    \item 口径 M（Model-usable）：在口径 I 基础上排除可计算型门控失败，用于标准 layout 指标（2D/3D IoU）。可比型门控失败仅影响对应指标。
\end{itemize}

\texttt{active\_time} 数据质量审计：在口径 T 中以审计表形式报告 \texttt{active\_time} 日志的异常记录比例，包括 \texttt{task\_id} = unknown、\texttt{annotator\_id} = unknown、\texttt{project\_id} = unknown、\texttt{script\_version} 缺失、以及多 session（task, annotator）对的占比，使数据质量可追溯而非仅口头声称。

\paragraph{Type 4（流程/格式失败）的前移控制与披露。}
对于可由 UI 约束的 Type 4 子类（多选元标签的空选与互斥冲突），本文采用提交时硬校验使其 ``尽量不可发生''，从而将问题从事后清洗前移为采集协议的一部分。为避免 ``挡在门外就当没发生'' 的粉饰风险，我们额外披露过程指标：被拦截/拒收的次数与原因分布（例如 \texttt{n\_rejected\_empty}、\texttt{n\_rejected\_conflict}、按 worker 的拒收率与重试次数分布）。
\emph{注意}：系统侧 NA（字段漏传、导出缺损、版本不一致）无法完全由 UI 消除，因此仍作为 Type 4/NA 的审计项在口径 T 中报告其发生率与分母。
\subsection{分层报告表}

本文通过四张分层表多角度展示结果：
\begin{itemize}
    \item 表 A：按 $\Delta n_{pairs}$ 分层——覆盖率、门控失败原因、改动幅度统计。
    \item 表 B：按 difficulty/\texttt{model\_issue} 分层——承认标签的主观性（每人主观感受不同），并报告 Fleiss' Kappa 验证一致性。
    \item 表 B2：混淆/冲突矩阵（低成本审计项）
Scope 冲突矩阵：In-scope vs OOS 及 OOS 子类冲突率
Tag 冲突统计：高频 difficulty/\texttt{model\_issue} 的共现与冲突比例；并强制披露两类\textbf{互斥冲突}与\textbf{空选}比例：
\begin{itemize}
    \item difficulty：\texttt{trivial} 与任意其他困难标签共存的比例；difficulty 空选/缺失比例；
    \item model\_issue（仅 semi）：\texttt{acceptable} 与任意其他 issue 共存的比例；model\_issue 空选/缺失比例。
\end{itemize}
用途：一方面作为分层解释与不确定性来源定位；另一方面作为\textbf{采集协议的可审计证据}：在提交时硬校验已启用的前提下，最终数据中的空选/互斥冲突应接近 0，同时应与 ``被拦截次数'' 的过程日志相互一致；若仍出现残余事件，则按任务列出示例并追溯原因（NA/导出缺损/版本不一致等）。
    \item 
表 C：Worker $\times$ Scene 可靠度矩阵
\\ 行：\texttt{worker\_id}
\\ 列：核心场景 $\mathcal{S}_{\mathrm{core}}$ 与 other。主文报告不超过 $S_{\mathrm{core,max}}=4$ 个核心场景，其余合并为 other 并单独披露频率与反例分布。
\\ 单元格：$r_{u}^{(s)}$ $\pm$ 95\% CI
\\ 标注：$N_{u,s}<N_{u,s,\min}$ 的格子标记为 --。本文保留 $N_{u,s,\min}\in\{5,6\}$ 两档备选，最终取值在 PreScreen 后冻结，并在附录~\ref{app:activation-threshold} 同时给出两档的保守启用率估计表。
\\ 覆盖率披露：报告各场景的有效样本分布、以及路由中 $r_u^{(s)}$ 实际启用比例（其余为退化到全局 $r_u$），以审计方式回答“场景分层是否在运行中退化为全局”的质疑。固定系数收缩版本仅在附录作为敏感性分析报告，不作为主实现。
\\ 补齐披露：报告 \texttt{Calibration\_reserve} 的实际使用量，并按用途拆分为 CI 追加与核心场景覆盖补齐两类。
\\ 降级披露：若总体启用率低于 $\tau_{act}$，则在主结论中不将场景路由作为强主张，仅将其作为可审计机制展示，并在结果中明确采用的降级口径。
\\ 类型标注：在每个 worker 行给出风险层级类型（$R0/R1/R2/R3$）及两轴坐标 $\mathrm{LCB}(r_u)$ 与离散风险层级位置，并在矩阵中标注其在不同场景与风险桶下的局部失效；同时附带主风险签名（blind-trust / condition-fragile / gate-failure）。
\\ 局部失效操作定义：预注册，满足以下任一条件即标记为 $failure_{u,s}=1$：
\begin{enumerate}
    \item 绝对阈值：$r_{u}^{(s)}<0.70$；
    \item 相对下降：$r_{u}^{(s)}<r_{u}-0.10$；
    \item 保守显著性：$r_{u}^{(s)}$ 的 95\% CI 上界 $<$ $r_{u}$ 的 95\% CI 下界。
\end{enumerate}

    \item 图 D：工人画像图
横轴：离散风险层级 $R_u^{\mathrm{tier}}\in\{R0,R1,R2,R3\}$
纵轴：可靠度 $\mathrm{LCB}(r_u)$
    \newline 说明：二维画像图是离散风险层级与连续可靠度主轴的\emph{可视化载体}，分组定义以方法章节中的顺序判定规则为准。
    \newline 颜色：4 类风险层级（Stable / Trust-vulnerable / Condition-fragile / Unusable-Noise）。
    \newline 形状：编码首要风险签名（blind-trust、condition-fragile、gate-failure、default stable），并与导出字段 \texttt{worker\_group\_reason} 保持一致。
    \newline 参考线：标注 $\tau_{r,low}$、$\tau_{r,high}$ 等可靠度阈值位置；横轴不再绘制连续异常度参考线，避免暗示伪精确的连续风险分数。
    \newline 可追溯性：报告侧对每位 worker 同时输出 \texttt{group\_rule\_version} 与 \texttt{worker\_group\_reason}，以支持复现实验与逐个核对。
\end{itemize}


\subsection{反例案例库}
输出内容：
\begin{itemize}
    \item 反例清单 CSV：task\_id, type, $\mathrm{IAA}_{t}$, $\mathrm{IoU}_{\mathrm{edit}}$, consensus\_tags, 原因
    \item  可视化：每个反例的标注对比图 3-4 个 worker 叠加
    \item 统计：三类反例的频率分布
    \item 人工复核笔记：专家确认的真正原因

\end{itemize}

 作用：必须说明其因果与证据链地位
 \begin{enumerate}
     \item 可审计筛选：依据第~\ref{sec:method-counterexample}节的可复现规则产生候选集，不依赖专家直觉选取样本。
     \item 专家复核：逐条确认是否为真正异常，并记录可解释原因与归类标签。
     \item 典型失败模式总结：形成可检索的失败模式词表，如跨门扩张、拓扑崩溃、配对失败、角点漂移等，并统计其在 IID 与 non-IID/OOD 分层下的频率变化。
     \item 反馈闭环：不等同于重训练闭环，将失败模式反馈到标注指南、培训素材、路由规则，例如在某类失败高发场景提高冗余或限制工人池。
     \item 过程性证据：在反例库中固定输出同一任务、不同功能组工人的对比可视化，用来解释为何预筛选与分组是必要的。
 \end{enumerate}

\subsection{过程性证据与归因链验证}\label{sec:report-process-evidence}
目的：将工人画像坐标、风险层级、Worker×Scene 矩阵、反例库与路由收益串联为可审计的证据链，而非仅报告最终 $p$ 值。
\\ 输出：
\begin{itemize}
    \item 同任务不同风险层级对比：选取若干典型任务，包含 OOD 与高风险，展示 $R0$ 与 $R1/R2$ 等不同风险层级的标注差异。
    \item 案例抽样规则：预注册，避免选择性展示，固定抽取 15 个任务用于正文展示，5 个高 $\mathrm{IAA}_{t}$ 共识明确，5 个低 $\mathrm{IAA}_{t}$ 分歧严重，5 个高 $d_{t}$ OOD 风险高；每个任务展示 3--4 名来自不同风险层级的标注叠加对比。
    \\ 与反例库关系：反例库候选集由第~\ref{sec:method-counterexample}节的可复现规则产生，其中低 $\mathrm{IAA}_{t}$ 组通常与 Type 1（低一致性）高度重叠；高 $d_t$ 组用于展示偏移下的典型行为，可能与 Type 2/3 重叠但不要求全部为反例。为保持可复现性，15 个任务的最终选取在进入 Main 前冻结排序规则与随机种子：各组先按对应指标排序，若并列则用固定随机种子打破并列；并在候选集中优先选择被反例规则标记的任务，否则从该组剩余任务中按固定随机种子等概率抽取。
    \item 失败模式溯源：对高频反例类型，回溯到对应场景与功能组，解释为何会发生、如何被预筛选或路由缓解。
    \item 路由收益归因：证明 stratified routing 的收益主要来自高风险任务分配给更稳定工人与 stress mode 增冗余或复核，而不是静默过滤。
\end{itemize}







\section{扩展：标注数据重训练（非主贡献）}

本节仅作为扩展讨论，不进入 RQ1--RQ3 的主结论。

若开展重训练，目标是利用反例案例库筛选更有信息量的数据，减少结构性错误（如跨门扩张与配对失败）并提高可计算性/可渲染性。门控失败率的变化仅作为安全性/数据可用性审计与诊断指标披露，不作为主收益主张来源。评估应固定测试集以避免数据泄漏，对比旧模型 $M_0$ 与新模型 $M_1$，并报告 \texttt{model\_issue} 频率、$\mathrm{IoU}_{\mathrm{edit}}$ 的变化，以及安全性事件（门控失败/字段缺失等）的比例与原因分布。

本文当前不做重训练的原因在于样本量与周期有限，且重训练会显著增加方案自由度与泄漏风险控制成本，削弱本文对规程、审计与路由证据链的聚焦。后续若扩展到更大规模且跨数据源的数据，可在保持与本文一致的预注册与隔离协议下再引入重训练闭环。



\section{讨论与局限性}

本章预先回应审稿人可能提出的质疑，对研究设计中的关键假设、方法边界与已知局限进行自我审查。
\subsection{Meta-label 的主观性}

Difficulty 与 \texttt{model\_issue} 依赖标注者主观判断，这是本研究的内在局限。
缓解措施包括统一培训并提供案例示范，通过 per-tag 共识降低个体偏差，并在 Calibration 上报告 Fleiss' Kappa 等一致性统计以量化噪声水平，计算口径见第~\ref{sec:method-tags-consensus}节。
本文的立场是：即使存在测量噪声，这些标签作为风险信号和场景代理仍具备实用价值。

\subsection{Type 4 的可控性边界（无法保证绝对为 0）}

从审稿视角，能将 ``字段缺失/不合规'' 从事后清洗前移到提交时强约束，是采集协议可控性的加分项。本文通过 UI 侧必填与互斥硬校验尽量使 ``空选'' 与 ``互斥冲突'' 在提交时不可发生，并同步记录被拦截次数作为过程证据。

但我们明确接受一个现实结论：Type 4 无法保证绝对为 0。即便 UI 侧硬阻断，系统仍可能出现 NA 或字段缺失（脚本未加载、客户端/服务端版本不一致、导出缺损、记录被污染等）。因此本文仍保留 Type 4/NA 的定义与审计披露，将其作为口径 T 的过程性质量指标，而不是在最终清洗后 ``消失''。
\subsection{样本量限制}
当前可用候选池为 458 张（\texttt{mp3d\_layout}/test）。在 10 人（标注者规模上限约 15--20 人）与紧周期约束下，主方案的阶段划分与规模见第~\ref{sec:exp-data-org}节，核心主分析使用约 260 张唯一图像（Calibration\_manual 100 + Validation 60 + Test 100）。
同理，标注者规模上限约 15--20 人时，细粒度画像分型容易出现空桶与阈值不稳定。为降低稀疏性与事后解释空间，主分析与主报告使用“可靠度主轴 + 离散风险层级”的 4 类工人画像，并将更细粒度的诊断型分型作为未来工作。
\\ 影响体现在：某些低频场景的 $r_{u}^{(s)}$ 仍可能因为样本数不足或置信区间过宽而不可用。低频场景的处理与退化规则见第~\ref{sec:method-scene-reliability}节；本文在结果中透明报告退化比例，并避免用不稳定的场景估计驱动激进路由。
\\ 补充说明：为避免在小样本下引入额外可调超参数而造成软门控调参质疑，主分析采用保守的 LCB 与退化策略，而不采用需额外超参数的收缩平滑或基于相关性的层级回退；相关扩展仅在未来工作讨论。
\\ 透明披露：在 Results 中报告每个场景的样本量和置信区间宽度。
\subsection{离线路由模拟的局限}
方法是在多标注数据上进行模拟分配后取该 worker 真实标注的离线路由评估。
局限在于这不属于真实在线路由，worker 行为可能因任务分配而改变。
辩护理由：离线模拟是路由研究的标准方法，参考 Kara et al. 和 Liu et al. 的工作；避免重新标注的巨大成本；为未来在线实验提供先验。离线模拟回答的是在固定已采集行为数据下策略的可解释收益下界；路由对标注者行为的正向在线效应如激励与适应性改变属于未来工作范畴，不混入主结论。
\subsection{Convex Hull 的局限性}
凸包外包凹形，可能降低对跨门扩张或凹结构错误的敏感度。本文仅将 $\mathrm{IoU}_{\mathrm{edit}}$ 作为改动幅度的辅助证据，其定义与边界已在第~\ref{sec:method-iou-edit}节说明；主结论优先依赖更几何一致的指标，并在反例筛选中采用多指标组合而非仅依赖 IoU。

\subsection{基本假设与让步}\label{sec:disc-assumptions}
受控组织假设：系统部署在受控组织内，如企业内部标注团队，主要面对噪声、平庸、不稳定而非强对抗的恶意攻击者。
\begin{itemize}
    \item 让步1：在极端 OOD 场景下可能出现多数标注者同时犯错而专家意见被淹没的情况。本文在高 $d_t$ 情形更依赖应激模式下的冗余提升与复核机制；加权共识的用法与边界见第~\ref{sec:method-weighted-consensus}节。
	\item 让步2：循环论证风险的规避，通过 PreScreen 的测试样本与专家参考标准获得 $r_{u}^{(0)}$ 的独立锚点，并由此映射得到 $w_u$，同时采用 leave-one-task-out 的估计协议，避免将同一任务的标签既作为推断依据又作为评价结论。
    \item 让步3：本文不在主分析中估计工人 OOD 韧性曲线这类需要在高 $d_{t}$ 区间大量覆盖的数据驱动函数，例如 $r_{u}(d_{t})$ 的衰减斜率，仅将 $d_{t}$ 作为触发器与分层变量；更精细的韧性建模留待更大规模与更长周期的后续研究。

\end{itemize}

\subsection{未来工作：细粒度工人画像分型}
本文在主分析中使用“可靠度主轴 + 离散风险层级”的 4 类工人画像（$R0/R1/R2/R3$）以保证小样本下的稳健性。若在更大规模的标注团队与更长周期下采集到更丰富的过程数据，可以在保持可靠度主轴不变的前提下，把当前风险层级进一步细化为更诊断化的行为类型，用于更精细的审计与培训反馈。

一种可行的细分方案，是在保持当前离散风险层级可追溯规则不变的前提下，再在每一级内部细分 6 类诊断型画像。其前提是：每类达到最小样本量，且 blind-trust、条件性脆弱与工程失败倾向的估计已经足够稳定，避免空桶导致的过拟合解释。
\begin{itemize}
	\item Reliable：高 $\mathrm{LCB}(r_u)$ 且未触发 blind-trust、条件脆弱或工程失败的核心风险。
	\item Normal：$\mathrm{LCB}(r_u)$ 中等，且仅有弱辅助审计风险，不触发主要风险层级升级。
	\item Trust-vulnerable：全局可靠度尚可，但在 semi 误导性初始化上表现出较高 blind-trust 倾向。
	\item Condition-fragile：全局可靠度尚可，但在高风险桶或特定场景上出现显著退化。
	\item Noisy/Sloppy：可靠度偏低并伴随较高工程失败倾向或持续性低质量过程行为。
	\item Severe unusable：在 manual anchors 或 calibration 中已无法满足最低 admission / main-pool 使用条件。
\end{itemize}
该细分的定位是诊断与解释，而非把小样本下的细类分布当作可泛化结论。主分析仍应优先依赖 $\mathrm{LCB}(r_u)$、离散风险层级与保守路由规则，并在高风险任务上采用更严格的候选池与冗余策略。



% 主线收敛到讨论与局限性；审稿问答与贡献总结暂不纳入主文。
% \section{与审稿员预期批评的预设回答}

以下针对审稿人可能提出的六个典型问题逐一给出回应。
问题1：Meta-label 会显著增加工作量吗？
回答：不会。Scope 是必要决策字段；Difficulty 与 \texttt{model\_issue} 属于修改时的同步记录。额外时间预计 $<10$ 秒，占平均标注时长约 8\%。
问题2：多选标签如何定义共识？
回答：采用 Per-tag 二元共识。每个 tag 独立多数票，报告 per-tag Fleiss' Kappa。为避免 ``空值=无困难'' 与 ``漏填'' 不可区分，采集协议在提交时对多选元标签执行硬校验：至少选 1 项（空选不可提交）且禁止两类互斥共选（\texttt{trivial}+其他困难标签；\texttt{acceptable}+其他失败模式）。NA 仍可能因系统/导出缺损出现，因此 NA 不计入有效人数并单独披露缺失率；同时披露被拦截/拒收次数与原因分布作为过程证据。
问题3：样本量有限，结论是否可靠？
回答：主终点预注册，使用 BCa bootstrap CI，报告效应量与 MDE，不以 $p$ 值作为唯一依据；同时披露各场景样本量与 CI 宽度。
问题4：离线路由模拟是否弱于在线实验？
回答：离线路由是当前阶段可复现且成本可控的标准做法，可用于形成在线实验先验；本文明确其边界并保留后续在线验证。
问题5：加权共识是否循环论证，OOD 下专家会否被淹没？
回答：权重锚点来自 PreScreen 测试样本与专家参考标准，并采用 leave-one-task-out，避免同任务自证；$w_{max}$ 在 PreScreen 阶段按预注册的方案 A（不足则方案 B）锁定后冻结，避免事后调参；高 $d_{t}$ 场景进入 stress mode，提高冗余或触发专家复核。
问题6：改进幅度是否过小？
回答：以预注册 MDE 判定工程意义；只有效应量跨过 MDE 且统计证据一致时，才声称改进成立。



% \section{预期贡献总结}

本文的预期贡献涵盖规程设计、系统设计、方法学与实证发现四个层面。
贡献 0：规程设计。提出 Pilot→PreScreen→Main 的三段式可审计规程，在进入主实验前预注册并锁定阈值、超参数与权重映射。流程允许在软件未完全自动化的情况下通过任务包分发执行，但必须保留可审计的样本隔离与分发清单。
贡献 1：系统设计。提出可审计的半自动标注体系。三级报告口径 T/I/M 避免选择偏差；所有门控失败与 OOS 可追溯；Meta-label 作为决策副产品提供分布监控信号。
贡献 2：方法学。预筛选测试样本、场景特异可靠度、OOD 触发路由的组合。$r_{u}^{(s)}$ 的估计协议与 LOO 原则；Per-tag 共识与分歧度计算；OOD 风险分数 $d_{t}$ 的预注册主实现与敏感性分析框架。
贡献 3：实证发现。在全景布局标注任务上验证，以最终实证结果为准。半自动初始化在相同图像对照下可降低 \texttt{active\_time}；场景与分布感知路由在相同预算下可提升一致性代理并改善预算效率；门控失败与字段缺失等作为安全性/数据可用性审计事件被完整追踪与披露（不作为主收益主张来源）；反例案例库揭示三类主要失效模式，并形成过程性证据链。
对标文献：Kara et al.（JAIR 2018）聚焦主动学习与共识估计，本文在此基础上引入场景分层；Liu et al.（InfoVis 2018）聚焦 worker 分型与交互验证，本文采用离线路由模拟加以拓展。创新点在于将可审计性、场景感知与元标签共识进行系统性整合。




\appendix
\section{附录：扰动算子库与误导性初始化冻结清单}\label{app:perturbation-library}

本附录给出 \texttt{PreScreen\_semi} 中误导性初始化的扰动算子库。算子库以现有 \texttt{model\_issue} 标签体系为索引，目标是生成可修复且可审计的结构性错误。Pilot 的作用是归档高频失效模式并锁定算子集合、强度范围与抽样配额。

为避免工程实现与论文披露的类型集合发生漂移，\texttt{model\_issue} 的取值集合与命名（alias）以 Label Studio 标注界面配置（\texttt{tools/label\_studio\_view\_config.xml}）为唯一真源；扰动算子实现与冻结清单中的 \texttt{operator\_id} 必须与该 alias 完全一致。

误导性初始化来源采取“双来源但主次分明”设计：主来源为规则化扰动算子，补充来源为固定小配额的自然失败 trap bank。自然失败来源仅允许进入 \texttt{PreScreen\_semi} 的盲信子集，不得与 \texttt{PreScreen\_manual} 复用同一 \texttt{base\_task\_id}。若某类型自然失败样本稀缺，则下调该来源的实际配额，并用同类型规则化扰动补足总 trap 配额；计划配额与实际配额需同时披露。

\subsection{扰动算子与 \texttt{model\_issue} 对应表}\label{app:perturbation-operators}

\begin{table}[ht]
\centering
\small
\begin{tabular}{p{0.18\linewidth} p{0.34\linewidth} p{0.18\linewidth} p{0.16\linewidth}}
\toprule
\textbf{\texttt{model\_issue}} & \textbf{扰动目标与算子示例} & \textbf{强度锁定} & \textbf{Pilot 观测}\\
\midrule
\texttt{acceptable} & 纠错子集使用高质量初始化 $P_t$，不施加强扰动，仅保留轻微扰动用于检验细调能力。 & 轻微 & 较多\\
\texttt{overextend\_adjacent} & 跨门扩张算子 $\mathcal{T}_{\mathrm{overextend}}$：将边界沿门洞方向外推并跨越门框停止点，制造可修复的跨门错误。 & 弱/中/强 & 中等\\
\texttt{corner\_drift} & 角点漂移算子 $\mathcal{T}_{\mathrm{drift}}$：对局部角点施加受限平移，保留拓扑结构但造成对齐偏差。 & 弱/中 & 较多\\
\texttt{corner\_duplicate} & 角点重复算子 $\mathcal{T}_{\mathrm{duplicate}}$：在同一物理拐角附近插入冗余角点，诱发配对不稳但仍可修复。 & 中/强 & 较多\\
\texttt{underextend} & 漏标算子 $\mathcal{T}_{\mathrm{underextend}}$：裁剪部分边界段或删除局部角点，使初始化覆盖不完整。 & 弱/中 & 极少\\
\texttt{over\_parsing} & 过度解析算子 $\mathcal{T}_{\mathrm{overparse}}$：插入伪角点或伪折线段，模拟把非结构细节误判为布局拐角的情形。 & 弱 & 极少\\
\texttt{topology\_failure} & 拓扑崩溃算子 $\mathcal{T}_{\mathrm{topology}}$：制造配对或闭合不稳定的拓扑错误。该模式在 Pilot 中罕见，因此不作为主流程的高占比来源，仅保留小配额用于稳健性披露。 & 固定小配额 & 罕见\\
\texttt{fail} & 预标注失效算子 $\mathcal{T}_{\mathrm{fail}}$：生成明显不可用但仍可从零重画的初始化。该模式在 Pilot 中极少，默认不进入盲信主指标集合，仅用于附录真实性检查。 & 固定小配额 & 极少\\
\bottomrule
\end{tabular}
\caption{扰动算子库与 \texttt{model\_issue} 的对应关系。强度为预注册锁定的离散档位，用于保证可复现并支持敏感性分析。Pilot 观测为 Pilot 样本内的相对频次描述，用于锁定抽样配额而非事后调整。}
\label{tab:perturbation-operator-library}
\end{table}

\subsection{冻结清单与可审计字段}\label{app:perturbation-freeze-list}

对每个误导性初始化任务，主实现冻结并记录以下字段，以保证从 $P_t$ 到 $\tilde{P}_t$ 的证据链完整。
\begin{itemize}
    \item \texttt{image\_id} 与任务唯一标识。
    \item \texttt{source\_type}（\texttt{synthetic\_operator} / \texttt{natural\_failure}）、\texttt{base\_task\_id} 与所属池标识，用于验证跨池不重叠。
    \item 模型版本标识与生成 $P_t$ 的配置摘要。
    \item 扰动算子标识 \texttt{operator\_id} 与强度档位 \texttt{lambda\_level}。
    \item 随机种子 \texttt{seed}。
    \item 由 $P_t$ 生成 $\tilde{P}_t$ 的脚本版本号与哈希。
\end{itemize}

自然失败初始化若被纳入补充来源，同样在固定候选池上运行冻结模型得到 $P_t$，再由专家按锁定的 \texttt{model\_issue} 归档表进行标签化，并按类型分层随机抽样进入盲信子集。该来源的图像清单、专家复核记录、抽样随机种子与计划/实际配额一并冻结。
\section{附录：数据集分组与工作量汇总}\label{app:dataset-summary}

\subsection{各数据集分组信息}\label{app:dataset-table}

表~\ref{tab:dataset-summary} 汇总了实验方案中所有数据集的唯一图像数、冗余设置、图像重叠关系与主要用途。
$N_{\mathrm{pool}}=458$，已分配唯一图像 355 张，保留 103 张用于补充与稳健性分析。

\begin{table}[H]
\centering
\caption{各阶段数据集分组汇总（$N_{\mathrm{pool}}=458$，已用唯一图像 355 张）}
\label{tab:dataset-summary}
\footnotesize
\resizebox{1.0\textwidth}{!}{%
\begin{tabular}{p{0.12\linewidth} p{0.18\linewidth} p{0.08\linewidth} p{0.06\linewidth} p{0.14\linewidth} p{0.36\linewidth}}
\toprule
\textbf{阶段} & \textbf{集合名称} & \textbf{条件} & \textbf{$N_{img}$} & \textbf{$k$ 配置} & \textbf{主要用途} \\
\midrule
Stage 0 & \texttt{Pilot} & manual & 15 & $k=2$ & 流程验证；归档 \texttt{model\_issue}；锁定扰动算子与强度范围 \\[3pt]
Stage 1 & \texttt{PreScreen\_manual} & manual & 30 & 全员 & 估计 $r_u^{(0)}$；其中 $n_{anchor}=12$ 为专家参考锚点；锁定 $w_{\max}$ \\[3pt]
Stage 1 & \texttt{PreScreen\_semi} & semi & 30 & 全员 & 测量纠错能力与盲信风险；与 \texttt{PreScreen\_manual} 图像池不重叠 \\[3pt]
Stage 2 & \texttt{Calibration\_manual} & manual & 100 & $k=4$ & 估计 $r_u$ 与 $r_u^{(s)}$（仅 manual）。内部划分为 \texttt{Calibration\_anchor} 12 张全员完成、\texttt{Calibration\_core} 75 张平衡分配、\texttt{Calibration\_reserve} 13 张用于追加 \\ [3pt]
Stage 2 & \texttt{Calibration\_semi} & semi & $(25)^{\dag}$ & $k=4$ & RQ2 配对对照（IAA/反例等）；从 \texttt{Calibration\_manual} 中预注册抽取子集；不用于估计 $r_u$ 或 $r_u^{(s)}$ \\ [3pt]
Stage 3 & \texttt{Manual\_Test} & manual & 100 & $k=1$ & RQ1 效率对比（手工侧） \\[3pt]
Stage 3 & \texttt{SemiAuto\_Test} & semi & $(100)^{\dag}$ & $k=1$ & RQ1 效率对比（半自动侧）；与 \texttt{Manual\_Test} 同图 \\[3pt]
Stage 3 & \texttt{Validation\_semi} & semi & 60 & $k_0=3,k_{\max}=5$ & RQ3 路由效果；OOD 分层；序贯停止 \\[3pt]
Stage 3 & \texttt{Gold\_manual} & manual & 20 & $k=3$ & RQ3 路由稳健性对照 \\
\midrule
\multicolumn{3}{l}{\textbf{合计（唯一图像，不重复计数）}} & \textbf{355} & & 余 103 张作为预留缓冲 \\
\bottomrule
\end{tabular}
}

\smallskip
\noindent\textit{注：括号内数值 $(n)^{\dag}$ 表示与同阶段另一集合共享同一批唯一图像，不另计入总数。}
\end{table}
\FloatBarrier

% \subsection{每位工人标注工作量估算}\label{app:workload-estimate}

% 以下估算假设 PreScreen 阶段全员完成两套图像池（各 30 张），Main 阶段按 per-image 冗余 $k$ 均匀分配。
% $\bar{k}_{\mathrm{Validation}}$ 取 $k_0=3$ 与 $k_{\max}=5$ 的均值 3.5 作为保守估计。

% \begin{table}[ht]
% \centering
% \caption{每位通过工人的预计标注任务数（上下界，$N=15$ 与 $N=20$）}
% \label{tab:workload-per-worker}
% \small
% \begin{tabular}{p{0.14\linewidth} p{0.22\linewidth} p{0.07\linewidth} p{0.07\linewidth} p{0.14\linewidth} p{0.14\linewidth}}
% \toprule
% \textbf{阶段} & \textbf{集合} & \textbf{$N_{img}$} & \textbf{$k$} & \textbf{$N=15$（每人）} & \textbf{$N=20$（每人）} \\
% \midrule
% Stage 1 & \texttt{PreScreen\_manual} & 30 & 全员 & 30 & 30 \\
% Stage 1 & \texttt{PreScreen\_semi}   & 30 & 全员 & 30 & 30 \\
% Stage 2 & \texttt{Calibration\_manual} & 30 & 4 & 8 & 6 \\
% Stage 2 & \texttt{Calibration\_semi}   & 30 & 4 & 8 & 6 \\
% Stage 3 & \texttt{Manual\_Test}        & 100 & 1 & 7 & 5 \\
% Stage 3 & \texttt{SemiAuto\_Test}      & 100 & 1 & 7 & 5 \\
% Stage 3 & \texttt{Validation\_semi}    & 60  & 3.5（均值）& 14 & 11 \\
% Stage 3 & \texttt{Gold\_manual}        & 20  & 3 & 4 & 3 \\
% \midrule
% \multicolumn{3}{l}{\textbf{手工小计（manual）}} & & \textbf{49} & \textbf{44} \\
% \multicolumn{3}{l}{\textbf{半自动小计（semi）}} & & \textbf{59} & \textbf{52} \\
% \multicolumn{3}{l}{\textbf{合计}} & & \textbf{108} & \textbf{96} \\
% \midrule
% \multicolumn{3}{l}{其中 PreScreen 占比} & & 60/108 $\approx$ 56\% & 60/96 $\approx$ 63\% \\
% \bottomrule
% \end{tabular}

% \smallskip
% \noindent\textit{注：PreScreen 由候选工人（约 25--30 人）完成，Main 仅由通过筛选的工人参与。
% 每个 PreScreen 图像池 30 张，其中 \texttt{PreScreen\_manual} 含 12 张专家参考锚点；每人 PreScreen 工作量总计 60（两个池），
% 使得总工作量降至约 108（$N=15$）或 96（$N=20$），PreScreen 占比降至约 56--63\%，实质性降低候选工人退出风险。}
% \end{table}

\section{附录：场景启用门槛的保守估计表}\label{app:activation-threshold}

本附录给出场景特异可靠度 $r_u^{(s)}$ 的启用门槛 $N_{u,s,\min}$ 在两种备选取值下的保守启用率估计表。门槛最终取值在 PreScreen 结束后冻结。

设核心场景集合为 $\mathcal{S}_{\mathrm{core}}$，其由 \texttt{Calibration\_anchor}+\texttt{Calibration\_core} 第一轮完成后按频率与一致性冻结得到。每位通过 PreScreen 的 worker 在 \texttt{Calibration\_manual} 至少获得 $N_{cal,\min}=25$ 个任务。对每个核心场景 $s$ 统计其频率并取保守下界 $\underline{p}_s$，用二项近似给出启用率下界：
\[\Pr\bigl(N_{u,s}\ge N_{u,s,\min}\bigr)\ge \Pr\bigl(\mathrm{Bin}(N_{cal,\min},\underline{p}_s)\ge N_{u,s,\min}\bigr).\]

\begin{table}[H]
\centering
\caption{启用门槛为 $N_{u,s,\min}=5$ 时的保守启用率估计表}
\label{tab:activation-nusmin-5}
\small
\begin{tabular}{p{0.18\linewidth} p{0.13\linewidth} p{0.13\linewidth} p{0.16\linewidth} p{0.18\linewidth}}
\toprule
\textbf{核心场景 $s$} & $\underline{p}_s$ & $N_{cal,\min}$ & $\mathbb{E}[N_{u,s}]$ & $\Pr\left(N_{u,s}\ge 5\right)$ \\
\midrule
$s_1$ &  & 25 &  &  \\
$s_2$ &  & 25 &  &  \\
$s_3$ &  & 25 &  &  \\
$s_4$ &  & 25 &  &  \\
\bottomrule
\end{tabular}
\end{table}
\FloatBarrier

\begin{table}[H]
\centering
\caption{启用门槛为 $N_{u,s,\min}=6$ 时的保守启用率估计表}
\label{tab:activation-nusmin-6}
\small
\begin{tabular}{p{0.18\linewidth} p{0.13\linewidth} p{0.13\linewidth} p{0.16\linewidth} p{0.18\linewidth}}
\toprule
\textbf{核心场景 $s$} & $\underline{p}_s$ & $N_{cal,\min}$ & $\mathbb{E}[N_{u,s}]$ & $\Pr\left(N_{u,s}\ge 6\right)$ \\
\midrule
$s_1$ &  & 25 &  &  \\
$s_2$ &  & 25 &  &  \\
$s_3$ &  & 25 &  &  \\
$s_4$ &  & 25 &  &  \\
\bottomrule
\end{tabular}
\end{table}

在主报告中同时给出两张表的数值，并披露实际启用比例与退化比例，用于审计场景分层是否在运行中退化为全局。

若第一轮校准后核心场景覆盖不足，则按第~\ref{sec:method-global-reliability}节的覆盖补齐规则，从 \texttt{Calibration\_reserve} 中对核心场景进行补派并再次报告启用率变化。若补齐后总体启用率仍低于 $\tau_{act}$，则主结论采用降级口径，不将场景路由作为强主张。
\section{附录：测度与统计方法速查}\label{app:metrics-stats-cheatsheet}

本附录以速查形式汇总本文的核心测度与统计学方法，目的为：在不改变主线研究问题（RQ1--RQ3）主终点定义的前提下，帮助审稿人快速理解「每个指标测什么、在哪个口径有效、如何统计检验、哪些属于主分析/敏感性分析」。

\subsection{测度总表}

\small
\begin{longtable}{p{3cm} p{2.5cm} p{2cm} p{1.5cm} p{3cm} p{3.5cm}}
\caption{测度总表（符号、层级、适用口径与用途）}\label{tab:metrics-overview}\\
\toprule
\textbf{测度} & \textbf{符号/字段} & \textbf{层级} & \textbf{口径} & \textbf{主要用途} & \textbf{说明/边界}
\\
\midrule
\endfirsthead

\multicolumn{6}{l}{\small（续表~\ref{tab:metrics-overview}）}\\
\toprule
\textbf{测度} & \textbf{符号/字段} & \textbf{层级} & \textbf{口径} & \textbf{主要用途} & \textbf{说明/边界}
\\
\midrule
\endhead

\midrule
\multicolumn{6}{r}{\small 续下页}\\
\endfoot

\bottomrule
\endlastfoot

活跃标注时间 & \texttt{active\_time} & task / worker & T / I & RQ1 主终点（效率）；可作为“成本维度”辅助汇总 & 仅度量交互活跃时间；需配套审计表披露 unknown/缺失与多 session；若用于综合分需披露覆盖率与异常率。\\
改动幅度（工作量代理） & $\mathrm{IoU}_{\mathrm{edit}}(t,u)$ & annotation & T / I & RQ1 风险控制；反例 Type 2 & 仅在展示初始化（semi）任务定义；不表示正确性。\\
一致性（任务内） & $\mathrm{IAA}_t$ & task & I & RQ2 主终点（质量/可靠性变化） & 定义为任务内 pairwise IoU 的中位数；在预注册配对子集上做配对对照。\\
LOO 一致性代理 & $\mathrm{IoU}(A_{t,u},C_t^{(-u)})$ & annotation & I & RQ3 指标（离线回放） & 以 LOO 共识 $C_t^{(-u)}$ 为参考；要求 $|U_t|\ge 3$，更稳健时建议 $\ge 4$。\\
结构差异信号（角点对数差） & $\Delta n_{pairs}(t,u)$ & annotation / task & T / I & 分层报告；反例/失败溯源 & 仅作诊断与分层变量；不作为硬门控或主终点。\\
边界 RMSE（几何差异） & $\mathrm{BoundaryRMSE}$ & annotation & I / M & 质量/反例诊断（可比型门控后） & 用边界曲线差异替代不稳定点匹配；若曲线不可构造则记为门控失败并披露。\\
角点匹配 RMSE（辅助） & $\mathrm{RMSE}_{\mathrm{match}}(t,u)$ & annotation & I / M & 质量诊断（可比型门控） & 仅在匹配稳定且覆盖率足够时启用；否则记为 NA 并披露缺失率。\\
全局可靠度（工人） & $r_u$（median） & worker & I & 路由排序；画像纵轴 & 仅用 \texttt{Calibration\_manual} 估计；报告 95\% CI 与 LCB，用于保守排序。\\
场景特异可靠度 & $r_u^{(s)}$（median） & worker $\times$ scene & I & 核心场景路由（RQ3） & 仅在核心场景 $|\mathcal{S}_{\mathrm{core}}|\le 4$ 启用；不足门槛则退化为 $r_u$ 并披露启用/退化率。\\
场景启用率/退化率 & \texttt{activation\_rate} / \texttt{degeneration\_rate} & scene / worker & I & 回答“场景分层是否静默失效” & 作为审计性披露项：不足门槛或 CI 过宽导致退化到全局 $r_u$。\\
门控失败率 & \texttt{gate\_fail\_rate} & task / worker & T / I / M & 安全性/数据可用性审计（非主终点） & 区分可计算型失败（影响口径 M）与可比型失败（仅影响对应列）；用于证明不静默过滤与失败原因可追溯，而非作为路由收益主证据。\\
预算效率（序贯冗余） & $k_{used}(t)$ / $\bar{k}_{used}$ & task / split & T / I & RQ3 指标（预算效率） & 与 $k_0,k_{\max}$ 配套报告；需披露达到 $k_{\max}$ 的比例（专家复核/单独审计队列）。\\
标注前分布风险（OOD 代理） & $d_t$ / $I_t^{\mathrm{OOD}}$ & task & T / I & stress mode 触发；分层报告 & 主实现：共享 pre-head latent $\rightarrow \mathrm{GAP}_W \rightarrow$ L2 归一化 $\rightarrow$ 第 $K$ 近邻距离；$I_t^{\mathrm{OOD}}=\mathbf{1}[d_t>\tau_d]$，其中 $\tau_d$ 由校准集 leave-one-out 分位预注册锁定；不用于静默过滤。\\
标注前结构风险（诊断触发） & $g_t$ & task & T / I & stress mode 触发；失败溯源 & 由初始化结构合法性/失败模式构造；不可计算时记 NA 并披露。\\
工人风险层级 & $R_u^{\mathrm{tier}}$ & worker & T / I & 画像横轴；审计/降级规则 & 离散风险层级（Stable / Trust-vulnerable / Condition-fragile / Unusable-Noise）；不使用单一连续异常分数作为主画像轴。\\
风险签名 & $T_u,C_u,G_u$（主）; $E_u,M_u$（辅） & worker & T / I & 支撑画像分组与路由解释 & 主风险分别对应盲信初始化、条件性脆弱、可计算失败倾向；低努力与缺失仅作辅助审计。\\
标准 Layout 指标（对齐基准） & 2D IoU / 3D IoU & task & M & 与 HoHoNet 等对齐的补充报告 & 仅在口径 M（可计算型门控后）报告；不用于样本静默剔除。\\
渲染深度指标（对齐基准） & depth RMSE / $\delta$ & task & M & 补充质量刻画（几何一致性） & 基于 layout 渲染得到深度图再比较；若渲染失败记为门控失败并披露。\\
OOS 比例与子类 & \texttt{scope=OOS} + 子类 & task & T & 数据可用性审计 & OOS 不进入主指标（I/M），但必须在 T 中报告比例、子类分布与冲突率。\\
多选标签共识/分歧 & $c_{t,k}$, $\mathrm{conf}_{t,k}$, $\mathrm{disagree}_t$ & task $\times$ tag & T / I & 场景代理；序贯追加诊断 & 仅用于分层、审计与追加触发（任务已获得足够元标签后）；不得用于静默过滤样本。\\
一致性统计（标签） & Fleiss' $\kappa$（per-tag） & tag / split & T / I & 证明 meta-label 可解释性 & 低频标签需单独标注 low-freq 警告；仅作审计性解释而非硬门控。\\

\end{longtable}
\FloatBarrier

\subsection{统计方法速查（主终点 + 辅助检验 + 敏感性分析）}

\begin{table}[H]
\centering
\caption{统计方法速查表（对应 RQ、主终点、检验与 CI）}
\label{tab:stats-overview}
\small
\begin{tabular}{p{0.10\linewidth} p{0.24\linewidth} p{0.20\linewidth} p{0.22\linewidth} p{0.20\linewidth}}
\toprule
\textbf{RQ} & \textbf{主终点/对比} & \textbf{效应量（主）} & \textbf{检验（主）} & \textbf{区间与备注}
\\
\midrule
RQ1 & Manual vs Semi 的 \texttt{active\_time} & Hodges--Lehmann 中位数差 或相对改进 & Mann--Whitney $U$（非参数） & 95\% BCa bootstrap CI（重采样次数预注册）；同时报告与 $\mathrm{IoU}_{\mathrm{edit}}$ 的关系作为风险控制。\\
RQ2 & 配对子集（$N_{img}=25$）的 $\mathrm{IAA}_t$：manual vs semi & $\mathrm{median}(\Delta_t)$，$\Delta_t=\mathrm{IAA}_t^{semi}-\mathrm{IAA}_t^{manual}$ & 同图配对置换检验（随机翻转符号） & 配对 bootstrap 95\% CI；K--S 仅附录敏感性分析，不作为主终点。\\
RQ2(反例) & 同配对子集的反例类型分布差异 & 类别比例差 / OR & $\chi^2$ 检验（或小样本用 Fisher） & 报告效应量与 CI（必要时 bootstrap）；与主终点一并计入多重检验校正。\\
RQ2(加权共识) & 加权 vs 未加权 $\kappa$/accuracy & $\Delta\kappa$ / $\Delta\mathrm{acc}$ & BCa bootstrap 重采样（任务级） & CI 不跨 0 且效应量 $\geq MDE_{\kappa}$ 认为显著；必要时用 McNemar 比较反例检出差异。\\
RQ3(离线回放) & 策略 3 vs 策略 1/2 的一致性代理、预算效率 & 任务级差异的中位数/均值差；$\bar{k}_{used}$ 差异 & 主分析优先配对置换/配对 bootstrap（t-test 仅附录敏感性） & 门控失败率与失败原因分布作为安全性审计披露（给出比例与 CI），不纳入主检验集合；若部署仅执行策略 3，则策略 1/2 仅在支持集做影子评估并披露支持率。\\
RQ3(部署) & $Validation_{semi}$ 上策略 3 的一致性代理（若可）、反例拦截、预算效率 & 差异/比值 & 配对置换/配对 bootstrap；二元率以 binomial CI 描述性披露 & 门控失败作为安全性事件仅报告比例、原因分布与上界风险，不作为“收益显著性”主张来源。\\
多重检验 & 计划检验集合（预注册） & --- & Bonferroni（保守） & 只对预注册检验集合校正；附录探索性分析单独标注，不与主结论混用。\\
\bottomrule
\end{tabular}
\end{table}

\subsection{方法实现细节速查（举例说明实现步骤）}

\begin{itemize}
	\item \textbf{加权共识的预注册实施：} 以 \texttt{PreScreen\_manual} 中 20--22 张 manual anchors 作为独立锚点，按第~\ref{sec:method-weighted-consensus} 节的方案 A（重复平衡二折交叉拟合作为主锁定方案）在每折内估计 $r_u^{(0)}$ 与初始 $w_u$，并在验证折上比较加权与未加权一致性（$\Delta\kappa$ 或 $\Delta\mathrm{acc}$）。3 折分层切分仅作附录敏感性分析；主分析不采用 5 折。选出跨重复切分均值最优且满足 $w_{max}\in\mathcal{W}$ 的值后即冻结，不再事后调参。若样本不足，使用方案 B 的 A/B 拆分方案，拆分种子与结果明细都需预注册并公开。
	\item \textbf{可靠度与场景特异性的 LOO 估计：} $r_u$ 与 $r_u^{(s)}$ 全部基于 \texttt{Calibration\_manual} 的 $|U_t|\geq3$ 样本，计算 LOO 共识 $C_t^{(-u)}$ 避免自我影响；若某场景启用次数不足（$|\mathcal{S}_{\mathrm{core}}|\leq4$），退化为全局 $r_u$ 并在附录中报告启用/退化率。置信区间由 bootstrap 获得，若半宽未达设定精度（预注册置信宽阈值），即时从 \texttt{Calibration\_reserve} 中追加任务。
	\item \textbf{序贯停止与预算效率计算：} Validation 中每张任务在候选池内按 $d_t$/$g_t$ 分层并触发应激模式，序贯增加 $k$ 直到 \texttt{stop\_threshold} 或 $k_{\max}$。平均每任务使用的 $\bar{k}_{used}$ 通过 paired bootstrap 比较；门控失败仅作为安全性审计事件记录其比例与原因分布，不纳入主检验集合。
	\item \textbf{Kolmogorov--Smirnov (K--S) 两样本检验（敏感性分析）：} 用于比较两组样本的分布（例如 $\mathrm{IAA}_t$ 在 manual vs semi 的分布形状）。计算经验分布函数 $F_n,G_m$ 并取最大差值 $D_{n,m}=\sup_x|F_n(x)-G_m(x)|$。实现上优先使用 \texttt{scipy.stats.ks\_2samp}（two-sided）；当存在大量并列值或样本量较小时，改用置换检验：在原样本上随机交换组标签 $B_{perm}$ 次（预注册建议 $B_{perm}=2000$ 或 $10000$）构造置换分布并计算双侧 p-value。报告原始 $D$ 值、p-value 与置换分布摘要，并在附录注明 K--S 对离散/并列数据的局限性。
	\item \textbf{Cohen's d（配对样本）：} 对配对连续量（如同任务下的均值差）计算配对效应量。令差分 $D_i=X_i^{(A)}-X_i^{(B)}$，配对 Cohen's d 定义为 $d_{paired}=\bar{D}/s_D$，其中 $s_D$ 为差分的样本标准差。实现要点：同时报告 $d$ 的 95\% CI，使用 paired bootstrap（预注册重采样次数 $B_{boot}=5000$ 或 10000）获得 CI；当差分分布严重偏离正态时，补报 Hodges--Lehmann 中位数差与其 bootstrap CI 作为稳健替代。解释阈值参考 small=0.2, medium=0.5, large=0.8，但以 MDE 与研究语境为准。
	\item \textbf{多重检验与 MDE 门槛：} 7 个预注册检验在 Bonferroni 下校正（$\alpha\approx 0.0071$），仅在 $p<\alpha$ 且效应量超过预注册 MDE（例如 Time $\geq8\%$ 改进、一致性差 $\geq0.05$、以及预算效率的最小可解释改进）时宣布显著；门控失败率不进入主检验集合，仅作安全性/可用性审计披露。MDE 取自 Pilot → PreScreen 的分级锁定流程，整个过程中不得事后调整阈值。
	\item \textbf{指标数据采集与一致性检验：} \texttt{active\_time} 以 session 分段追踪（见第~\ref{sec:method-active-time} 节），每种汇总（per session/per annotator/per task）都按最长累计值汇总，并通过 Spearman 相关与原生 \texttt{lead\_time} 比值的中位数、IQR 验证采集逻辑一致性。
\end{itemize}

\subsection{关于加权几何平均综合分的预注册使用边界}

若需报告综合分以帮助读者快速把握整体趋势（类比综合跑分），建议将其定位为\textbf{辅助汇总指标}而非替代 RQ1--RQ3 主终点。为降低调参空间，权重与归一化方式需在实验前冻结，且不得用于选择样本、选择策略或调参。

\paragraph{审稿友好的“最小自由度”建议（权重与敏感性分析）}
\textbf{(1) 只对同一研究问题内的输出做汇总：}优先将综合分限定在 RQ3（路由策略对比）的输出维度内（例如一致性代理与预算效率；若该集合确有可靠的时间日志，可额外加入时间维度）。门控失败率作为安全性/可用性审计项单独披露，不建议纳入综合分（稀有事件在小样本下会导致不稳定权重争议）。不建议跨 RQ1--RQ3 把时间、质量、路由收益混成单一总分，避免因口径与缺失机制不同引发选择偏差争议。

	\textbf{(2) 权重只给少量离散方案并全部报告：}避免连续权重“网格搜索最优”。建议预注册并固定报告少量离散权重（示例，按 2 维：一致性代理/预算效率）：
\begin{itemize}
	\item 等权：$w=(0.50,\,0.50)$（默认；最不易被指控偏好性）。
	\item 质量优先：$w=(0.67,\,0.33)$（强调一致性代理）。
	\item 成本优先：$w=(0.33,\,0.67)$（强调预算效率）。
\end{itemize}
\textbf{若将 \texttt{active\_time} 纳入跑分：}建议仅在“用于综合分的任务集合”上满足以下条件时启用：时间字段覆盖率 $\ge 90\%$、unknown/缺失异常率低于预注册阈值、且通过 \texttt{active\_time} 与 \texttt{lead\_time} 的一致性审计（见主文 RQ1 的一致性校验）。此时将时间作为第 4 个“成本维度”（越小越好，进入综合分前需做单调变换为越大越好），并同样仅报告离散的 3 套权重，例如：
\begin{itemize}
	\item 等权3维：$w=(1/3,\,1/3,\,1/3)$（质量/预算/时间）。
	\item 质量优先3维：$w=(0.50,\,0.25,\,0.25)$。
	\item 成本优先3维：$w=(0.25,\,0.50,\,0.25)$。
\end{itemize}
其中时间维度的单调变换需预注册锁死（例如基于 baseline 的相对改进并做截断），并强制报告时间维度的覆盖率与异常率，避免因缺失机制引发样本切换。

\textbf{(3) 只做低成本稳健性披露：}附录报告 leave-one-dimension-out（每次去掉一个维度重算综合分）以及几何平均 vs 算术平均的替代聚合，检查结论方向与排序是否一致；并强制披露综合分的有效样本覆盖率（complete-case vs partial-case），避免因门控/缺失导致的隐性样本切换。
% Ensure all appendix floats (tables/figures) are flushed before bibliography
\clearpage
\nocite{kara2018,liu2018}
\bibliographystyle{IEEEtran}
\bibliography{refs/references}

\end{document}
